\section{Results}
\label{sec:results}

\revblue{Results are reported first for Gazebo simulation (single and multiple stimuli, and gait switching on the three-topology terrain map described earlier) and then for the physical platform.} 

\subsection{Simulation Results}

In the first phase, the robot's behavior was analyzed under a single type of stimulus, following an approach similar to that presented by \citet{sarvestani2013computational}. Then, to assess performance in more complex environments, the platform was exposed to multiple stimuli of different modalities. Simulations were performed to verify that the selected gait corresponded to the most appropriate mode based on the perceived terrain. This stage emphasized the analysis of neuronal responses in the obstacle sensory module, the locomotor decision network, and the locomotion-mode decision network.

Additionally, to contextualize the simulation within its intended inspection application domain, a Gazebo environment inspired by the spatial layout of an industrial warehouse was constructed (Figure~\ref{fig:depot_simulation}); its scope and caveats as a symbolic representation rather than a certified test arena are discussed in Section~\ref{ssec:Limitations}.

\begin{figure}[t]
    \centering
    \includegraphics[width=0.75\textwidth]{Depot_simulation.jpeg}
    \caption{%
        Gazebo simulation environment used for stimulus-response
        experiments, inspired by the general spatial organization of an
        industrial storage facility (open corridors, distributed obstacle
        elements, defined navigation zones). See Section~\ref{ssec:Limitations}
        for the scope and caveats of this representation.%
    }
    \label{fig:depot_simulation}
\end{figure}

\subsubsection{Single Stimulus Response}

Figure~\ref{fig:RESPUESTAS_EST_UNI} shows the simulated trajectories for the three stimulus types, sampled at six fractions of the total trial time: (a) approach toward the appetitive stimulus, (b) evasive behavior in the presence of an obstacle, and (c) escape response to an aversive stimulus. The evasive and escape behaviors differ mainly in the locomotion mode employed and the speed of orientation change.

\begin{figure}[t]
    \centering
    \includegraphics[width=0.8\linewidth]{RES_EST_UNI.png}
     \caption{Response to single stimulus in simulation trials. \revblue{Each panel shows a single trial ($N=1$ per panel).}}
     \label{fig:RESPUESTAS_EST_UNI}
\end{figure}

These simulated behaviors arise from the neuronal dynamics in Figure~\ref{fig:NEU_EST_UNIG}: with an obstacle on the robot's right side, the platform evades left, driven by the sensory ring, whose unit $L_{4}$ activity progressively decreases as the robot moves away.

\begin{figure}[t]
    \centering
    \includegraphics[width=1.0\textwidth]{NEU_EST_UNIG.png}
    \caption{Raster plot of neural activations for single obstacle stimulus in simulation trial. \revblue{Single trial ($N=1$).}}
    \label{fig:NEU_EST_UNIG}
\end{figure}

Figure~\ref{fig:NEU_EST_UNIG}b shows the locomotor decision network: the obstacle-associated GPe activation triggers unit $X_{4}$, but inhibitory neuromodulation from $X_{11}$/$X_{12}$ limits the sensory/integrative layers' influence, so the lateral command is instead driven by $L_{0}$ and $L_{3}$ (front-right quadrant). Neuron $X_{17}$ stays inactive (no approach behavior), and the locomotion-mode module activates unit $X_{14}$, enabling omnidirectional gait.

For the appetitive scenario (target initially to the robot's right), Figure~\ref{fig:NEU_EST_UNIB} shows the locomotor decision network driving the approach: neuron $X_{6}$ generates the rightward centering movement, then $X_{13}$ drives the forward approach, and $X_{17}$ issues the stop signal near the inspection area; the locomotion-mode module activates unit $X_{16}$, enabling differential rolling mode.

\begin{figure}[t]
    \centering
    \includegraphics[width=1.0\textwidth]{NEU_EST_UNIB.png}
    \caption{Raster plot of neural activations for single appetitive stimulus in simulation trial. \revblue{Single trial ($N=1$).}}
    \label{fig:NEU_EST_UNIB}
\end{figure}

\subsubsection{Quantitative Validation: Single-Stimuli Simulation Performance}
The performance of the platform was evaluated under isolated sensory stimuli to validate the locomotor decision network and the sensory modules individually.

\begin{figure*}[t]
    \centering
    \includegraphics[width=\textwidth]{fig1_noise_robustness_grid.png}
    \caption{Combined perception-decision noise-robustness battery (Section~\ref{ssec:NoiseRobustnessEval}, Table~\ref{table:PerturbationParams}) across the four single-stimulus/multi-stimulus simulation scenarios (rows: Appetitive, Aversive, Obstacle, Complex; three retained trials/severity level/scenario) against four metrics (columns: Mission Time $T_{sim}$, Pitch RMS, Decision Latency, $T_{switch}$), all as a function of the Combined Perturbation Level (index $0$--$4$). \revblue{Experimental limitation: at every non-zero level, IMU drift, LiDAR range noise, and camera illumination distortion are injected jointly, not in isolation, so this battery cannot separate or rank each sensor's individual contribution to the observed robustness; in addition, the IMU component reflects cumulative random-walk drift only, not additive white noise.} Error bars denote $\pm 1\sigma$ across the three trials retained per point. The Pitch RMS column uses a per-row scale (ceiling at $1.1\times$ that row's maximum) rather than a fixed $0$--$4^\circ$ range, so that rows with little variation remain readable without losing cross-scenario magnitude comparability.}
    \label{fig:noise_robustness_grid}
\end{figure*}

\begin{figure}[t]
    \centering
    \includegraphics[width=0.75\textwidth]{fig2_Appetitive_micro_test_001_SUCCESS.png}
    \caption{Neural and kinematic response during a representative Appetitive Targeting trial at the unperturbed baseline (Combined Perturbation Level 0). \revblue{The trial shown is the first by index of the suite's $N=15$ trials (three per perturbation level), not selected by outcome.}}
    \label{fig:appetitive_micro}
\end{figure}

\textbf{Appetitive Targeting (see Figure~\ref{fig:noise_robustness_grid}):} The appetitive module retains task-level resilience across the full combined-perturbation sweep (Section~\ref{ssec:NoiseRobustnessEval}): every retained trial reaches the target regardless of severity level (see Figure~\ref{fig:noise_robustness_grid} for the per-level $T_{sim}$/Pitch RMS/Decision Latency/$T_{switch}$ trends). \revblue{This robustness is evaluated under the combined perturbation --- IMU drift, LiDAR noise, and illumination distortion injected jointly at every non-zero severity level --- rather than against any single sensor's contribution in isolation, so the trends in Figure~\ref{fig:noise_robustness_grid}, across all four scenarios, cannot be attributed to any single sensing modality. Moreover, the IMU channel is perturbed only through cumulative random-walk drift: additive IMU white noise, although implemented in the simulated sensor model, was held at zero across all five severity levels (Section~\ref{ssec:NoiseRobustnessEval}).}

\begin{figure}[t]
    \centering
    \includegraphics[width=0.75\textwidth]{fig2_Aversive_micro_test_001_SUCCESS.png}
    \caption{Trajectory deflection overdrive phase for a representative Aversive Escape trial at the unperturbed baseline (Combined Perturbation Level 0), showing the immediate computational response of the repulsive network blocks. \revblue{The trial shown is the first by index of the suite's $N=15$ trials (three per perturbation level), not selected by outcome.}}
    \label{fig:aversive_micro}
\end{figure}

\textbf{Aversive Escape (see Figure~\ref{fig:aversive_micro}):} The defensive repulsion properties of the architecture were re-evaluated under the combined-perturbation sweep, with every retained trial across all five severity levels resolving the threat successfully (Figure~\ref{fig:noise_robustness_grid}). The neural firing variance exhibits a swift, step-like onset upon threat detection — structural mobilization into a defensive overdrive state — and returns to baseline once the threat leaves the field of view.

\begin{figure}[t]
    \centering
    \includegraphics[width=0.75\textwidth]{fig2_Obstacle_micro_test_001_SUCCESS.png}
    \caption{Time-series synchronization for a representative Obstacle Evasion trial at the unperturbed baseline (Combined Perturbation Level 0), showcasing the sensory-locomotor coupling triggered by LiDAR proximity boundaries. \revblue{The trial shown is the first by index of the suite's $N=15$ trials (three per perturbation level), not selected by outcome.}}
    \label{fig:obstacle_micro}
\end{figure}

\textbf{Obstacle Evasion (see Figure~\ref{fig:obstacle_micro}):} Smooth capsule geometries for simulated obstacles prevent localized torque spikes during contact. Every retained trial across all five combined-perturbation levels completed evasion successfully (Figure~\ref{fig:noise_robustness_grid}). Neural firing variance stays at a low baseline until the robot enters the obstacle's ranging threshold, at which point the avoidance sub-circuits modulate steering weights while preserving forward momentum.

\subsubsection{Response to Multiple Stimuli of Different Natures}

For this experiment, an obstacle was placed close on the robot's right side, an aversive stimulus on the left, and an appetitive stimulus farther away, toward which the platform was expected to approach. Neuronal activities of the four control structures are shown in Figure~\ref{fig:NEU_EST_MUL}, and the resulting trajectory and behaviors in Figure~\ref{fig:RES_EST_MUL}.

\begin{figure}[t]
    \centering
    \includegraphics[width=\textwidth]{NEU_EST_MUL.png}
    \caption{Raster plot of the neuronal activations to multiple stimuli of diverse nature in a simulation trial. \revblue{Single trial ($N=1$).}}
    \label{fig:NEU_EST_MUL}
\end{figure}

\begin{figure}[t]
    \centering
    \includegraphics[width=0.8\textwidth]{RES_EST_MUL.png}
    \caption{Trajectory followed in response to multiple stimuli of different natures in a simulation trial. Snapshots captured at selected time instants (0, 8, 15, 20, 25, and every 35~s thereafter until 3~min). \revblue{Single trial ($N=1$).}}
    \label{fig:RES_EST_MUL}
\end{figure}

At the initial moment, the evasive sensory structure (Figure~\ref{fig:NEU_EST_MUL}b) and the obstacle-associated globus pallidus activate, signaling an evasive maneuver; in the locomotion structure (Figure~\ref{fig:NEU_EST_MUL}c), units $L_0$, $L_1$, and $L_3$ activate, indicating the obstacle is on the right. As it is cleared, activity in both modules decreases.

At $t=21$~s, the basal ganglia module (Figure~\ref{fig:NEU_EST_MUL}a) shifts dominance to the closer aversive stimulus: the robot flees rightward, driven by units $X_4$/$X_5$, with the evasion direction set by $X_{11}$ outweighing $X_{12}$; as with the obstacle, the aversive globus pallidus activation then declines, marking the behavioral transition.

At $t=27$~s the appetitive stimulus enters view, activating its globus pallidus and triggering approach: units $X_6$/$X_3$ integrate into $X_{13}$'s rightward turn, after which the robot advances straight until $X_{17}$'s stop signal.

The platform reaches the goal in 92~s, sequentially adjusting locomotion mode via units $X_{14}$ (obstacle evasion, omnidirectional), $X_{15}$ (aversive escape, quadrupedal), and $X_{16}$ (appetitive approach, differential rolling). In Figure~\ref{fig:NEU_EST_MUL}d, units $X_1$/$X_2$ stay inactive (no terrain instability detected, gait unchanged), while $X_0$ shows only a small-amplitude activation producing transient states around the quadrupedal selection, attributable to that gait's abrupt movements.

\subsubsection{Quantitative Validation: Multi-Modal Simulation Performance}

\begin{figure}[t]
    \centering
    \includegraphics[width=0.75\textwidth]{fig2_Complex_micro_test_001_SUCCESS.png}
    \caption{Temporal dynamics for a representative Complex Navigation trial at the unperturbed baseline (Combined Perturbation Level 0), demonstrating decision-making under conflicting sensory modalities. \revblue{The trial shown is the first by index of the suite's $N=15$ trials (three per perturbation level), not selected by outcome.}}
    \label{fig:complex_micro}
\end{figure}

\textbf{Complex Navigation (see Figure~\ref{fig:noise_robustness_grid}):} The complex navigation suite presents the obstacle, aversive, and appetitive stimuli within a single trial. In the neural firing-variance profile of Figure~\ref{fig:complex_micro}, the robot first clears the physical obstacle through the LiDAR-driven evasion pathway, without a corresponding shaded region since this channel is not one of the two stimulus-presence traces plotted; the aversive threat then drives a red-shaded burst of firing-variance activity as the quadrupedal escape response engages; once clear of the threat, a blue-shaded interval marks sustained appetitive tracking as the platform reorients and approaches the target under differential-rolling mode. The kinematic panel's Roll/Pitch RMS response tracks this same three-stimulus sequence, rising through each gait transition and settling once the target is reached. Every retained trial across all five combined-perturbation severity levels resolves this three-stimulus encounter successfully (Figure~\ref{fig:noise_robustness_grid}), evidencing that the arbitration network does not suffer from computational deadlocks even when a single trial exercises all three behavioral channels under joint sensor perturbation.

Table~\ref{tab:consolidated_simulation_metrics} consolidates completion-time and attitude statistics across all four scenarios; per-severity-level trends of these same metric families under the combined IMU-drift/LiDAR-noise/illumination-distortion perturbation (Section~\ref{ssec:NoiseRobustnessEval}) are reported separately in Figure~\ref{fig:noise_robustness_grid}. Values below are recomputed directly from the archived per-trial simulation logs (one Obstacle trial with an anomalous, physically implausible roll reading was excluded, as in Table~\ref{tab:fsm_comparison}).

\begin{table}[htbp]
\centering
\caption{Consolidated Quantitative Performance Metrics.}
\label{tab:consolidated_simulation_metrics}
% Definimos columnas centradas en X
\begin{tabularx}{\textwidth}{l >{\centering\arraybackslash}X >{\centering\arraybackslash}X >{\centering\arraybackslash}X >{\centering\arraybackslash}X}
\toprule
\textbf{Suite} & \textbf{Success (\%)} & \textbf{$T_{sim}$ (s)} & \textbf{Roll (deg)} & \textbf{Pitch (deg)} \\
\midrule
Appetitive Targeting & 100.0\% & 38.52 $\pm$ 1.20 & 0.54 $\pm$ 0.01 & 1.07 $\pm$ 0.00 \\
Aversive Escape      & 100.0\% & 9.98 $\pm$ 2.07  & 1.36 $\pm$ 0.11 & 1.11 $\pm$ 0.07 \\
Obstacle Evasion     & 100.0\% & 15.72 $\pm$ 2.49 & 1.55 $\pm$ 0.15 & 0.90 $\pm$ 0.31 \\
Complex Navigation   & 100.0\% & 24.96 $\pm$ 10.36 & 1.03 $\pm$ 0.35 & 1.01 $\pm$ 0.13 \\
\bottomrule
\end{tabularx}
\end{table}

\subsubsection{Combined-Perturbation Robustness: Terrain Stability, Failure Modes, and Camera Dropout}
\label{sssec:noise_robustness_extra}

Two further scenarios --- the Rugged Terrain course and the Inclined Slope, already introduced in Section~\ref{sssec:variable_topo_quant} --- were re-evaluated under the same combined IMU-drift/LiDAR-noise/illumination-distortion sweep as the four scenarios above (Section~\ref{ssec:NoiseRobustnessEval})\revblue{, with results reported in Figure~\ref{fig:noise_robustness_extra_combined}. As in those four scenarios, the three perturbation sources are injected jointly at every non-zero severity level rather than in isolation, so the results discussed below cannot be decomposed into each sensor's individual contribution; the IMU channel, as before, is perturbed only through cumulative random-walk drift, with additive white noise held at zero throughout (Section~\ref{ssec:NoiseRobustnessEval}).} This is a separate battery from the Rugged Terrain/Inclined Slope trials analyzed in Section~\ref{sssec:variable_topo_quant}\revblue{, which carried only proportional Gaussian sensor noise,} and is not directly comparable to it; the two batteries differ in perturbation condition, not just sample size.

\begin{figure*}[t]
    \centering
    \begin{subfigure}[t]{0.32\textwidth}
        \centering
        \includegraphics[width=\textwidth]{fig_R1_failure_by_family.png}
        \caption{Trial outcome by scenario.}
        \label{fig:failure_by_family}
    \end{subfigure}
    \hfill
    \begin{subfigure}[t]{0.32\textwidth}
        \centering
        \includegraphics[width=\textwidth]{fig_R2_terrain_stability_vs_noise.png}
        \caption{Postural stability vs.\ perturbation level.}
        \label{fig:terrain_stability_vs_noise}
    \end{subfigure}
    \hfill
    \begin{subfigure}[t]{0.32\textwidth}
        \centering
        \includegraphics[width=\textwidth]{fig_R3_camera_dropout_evidence.png}
        \caption{Camera-dropout evidence.}
        \label{fig:camera_dropout_evidence}
    \end{subfigure}
    \caption{Combined-perturbation robustness battery\revblue{ (Section~\ref{ssec:NoiseRobustnessEval}). Experimental limitation: as in Figure~\ref{fig:noise_robustness_grid}, IMU drift, LiDAR range noise, and camera illumination distortion are injected jointly at every non-zero severity level, which precludes isolating each sensor's individual contribution to the results shown in panels (a)--(c); the IMU channel reflects cumulative random-walk drift only, with additive white noise held at zero.} (a) Trial outcome by scenario across all six suites, all severity levels pooled; each bar decomposes total attempts into successful trials retained for analysis and discarded failures by cause (timeout, tipover, crash) --- the Rugged Terrain suite accounts for the large majority of non-timeout failures (8 tipovers, 4 crashes out of 27 attempts), markedly more than any other suite. (b) Chassis Roll RMS and Pitch RMS versus Combined Perturbation Level (index $0$--$4$) for the Rugged Terrain and Inclined Slope suites, successful trials only; Roll RMS on Rugged Terrain remains within a $1.65^\circ$--$1.9^\circ$ band across all severity levels, Pitch RMS on the Inclined Slope ranges $2.3^\circ$--$5.7^\circ$ across the sweep (the perturbation index is left as an unconverted integer rather than a percentage, since no single percentage scale represents all three sensor modalities' differently-defined severity units). (c) Appetitive-stimulus detection fraction over time for a matched pair of Appetitive Targeting attempts under camera illumination distortion: a failed attempt (timeout) and its successful retry at the same severity level --- the failed attempt's detection fraction collapses to $\approx 0.01$ for the majority of the trial, while the successful retry sustains $\approx 0.99$ detection throughout.}
    \label{fig:noise_robustness_extra_combined}
\end{figure*}

Figure~\ref{fig:failure_by_family} shows that failure under combined perturbation is not evenly distributed across scenarios: the Aversive Escape and Obstacle Evasion suites completed every attempt without a single discarded trial, the Appetitive Targeting and Complex Navigation suites discarded a share of attempts to timeout (respectively, before three successful replicates per severity level were retained), and the Rugged Terrain suite alone accounts for 8 tipovers and 4 crashes out of 27 attempts --- the only suite where the robot's structural stability, not merely task completion time, was compromised under perturbation. This asymmetry is addressed explicitly in Section~\ref{ssec:Limitations}.

Among the trials that complete successfully, Figure~\ref{fig:terrain_stability_vs_noise} shows postural stability does not visibly degrade with perturbation severity on either terrain: Roll RMS on Rugged Terrain and Pitch RMS on the Inclined Slope both stay within the bands reported in the caption across the full severity sweep, without a monotonic trend toward instability. Read together with Figure~\ref{fig:failure_by_family}, this indicates that Rugged Terrain's elevated tipover/crash count is concentrated in the trials discarded before completion rather than reflected as progressive attitude degradation in the trials that do complete --- a distinction relevant to how this suite's robustness should be characterized (Section~\ref{ssec:Limitations}).

Figure~\ref{fig:camera_dropout_evidence} isolates a case where a purely sensor-level event --- the illumination-distortion model's probabilistic detection dropout suppressing the appetitive-stimulus detector almost entirely for one attempt --- propagates all the way to a task-level failure: with the target undetected for nearly the whole trial, the appetitive channel never accumulates enough sustained activation to drive an approach, and the trial times out. The retry at the identical severity level, drawing a different dropout realization, sustains detection throughout and completes normally. This is precisely the pipeline-level (not sensor-level-only) evidence the combined-perturbation evaluation is designed to surface: the failure mode originates in the camera-illumination stage but only becomes visible as a robustness concern once its consequence for task completion is measured, not from the sensor-level distortion metric alone.

Regarding the LiDAR-noise component of this evaluation specifically: the perturbation reaches the reactive obstacle-evasion pathway (the sensory-ring and locomotion-module units described in Section~\ref{sec:methodology}) in every scenario reported above, but --- in the code that generated this dataset --- does not reach the basal-ganglia arbitration competition between the obstacle, appetitive, and aversive channels, whose LiDAR-derived input channel is held at zero irrespective of sensor noise. The robustness demonstrated here for LiDAR noise should accordingly be read as established at the reactive-evasion level, not as a demonstration that the three-way arbitration competition itself tolerates LiDAR noise.

\FloatBarrier

\subsubsection{Basal Ganglia Ablation Study}

To directly address the reviewer's request for an ablation analysis of the basal-ganglia arbitration network and a quantitative comparison against simple threshold-based decision logic, the four-variant ablation described in Section~\ref{sec:methodology} was evaluated across the Appetitive Targeting (no-conflict) and Complex Navigation (three-stimulus conflict) scenarios, reusing the same scenario parameters as Table~\ref{tab:consolidated_simulation_metrics} ($n=12$ trials/variant/scenario, noise level fixed at zero to isolate the structural effect of each ablation from sensory-noise robustness, which is characterized separately, earlier in this section, for the intact circuit). This study is simulation-only; no equivalent ablation was run on the physical platform.

Without stimulus conflict, all four variants perform statistically indistinguishably: $91.7$--$100\%$ success in the Appetitive Targeting scenario, $p=1.000$ for all three comparisons against the intact circuit --- removing an ablation target does not break the simple case, as expected when only one channel is ever active. Under real conflict (Complex Navigation, Figure~\ref{fig:ablation_success_rate}), a clear gradient emerges. The two ablations that remove basal-ganglia mechanism itself --- \emph{no lateral inhibition} (inhibitory connections) and \emph{no STN/STR} (neural layers) --- collapse the success rate to $25.0\%$ and $33.3\%$ respectively (Fisher's exact $p<0.001$ and $p=0.001$ vs.\ the intact circuit's $100\%$), directly addressing the reviewer's request to quantify the consequence of removing inhibitory connections or key neural layers (R3-01). Replacing the entire circuit with static threshold logic (\emph{threshold only}) degrades success more mildly and not significantly on this metric alone ($83.3\%$, $p=0.478$); the case against threshold logic instead rests on the continuous metrics below.

\begin{figure*}[t]
    \centering
    \begin{subfigure}[b]{0.42\textwidth}
        \centering
        \includegraphics[width=\linewidth]{fig_ablation_success_rate.png}
        \caption{Success rate\revblue{ under three-stimulus conflict ($n=12$ trials/variant; Fisher's exact $p$ vs.\ the intact circuit).}}
        \label{fig:ablation_success_rate}
    \end{subfigure}
    \hfill
    \begin{subfigure}[b]{0.55\textwidth}
        \centering
        \includegraphics[width=\linewidth]{fig_ablation_complex_stability.png}
        \caption{Attitude stability (Roll/Pitch RMS)\revblue{: mean per variant, error bars show the 95\% confidence interval ($n=12$ trials/variant; Mann--Whitney $p$ vs.\ the intact circuit).}}
        \label{fig:ablation_complex_stability}
    \end{subfigure}
    \caption{Basal-ganglia ablation study under real stimulus conflict (Complex Navigation scenario, $n=12$ trials/variant; Fisher's exact / Mann--Whitney $p$ vs.\ the intact circuit). \revblue{The four variants are: \emph{intact circuit} (control); \emph{no lateral inhibition} (cross-channel inhibitory terms zeroed); \emph{no STN/STR} (STN and STR layers removed, stimulus routed directly into GPe/GPi); and \emph{threshold only} (circuit replaced by a static fixed-priority threshold over the raw stimulus).} (a) Success rate: the two ablations that remove basal-ganglia mechanism itself collapse success rate significantly (\emph{no lateral inhibition}: $25.0\%$, $p<0.001$; \emph{no STN/STR}: $33.3\%$, $p=0.001$); replacing the circuit with static threshold logic degrades success more mildly and not significantly at this sample size (\emph{threshold only}: $83.3\%$, $p=0.478$). Under the no-conflict Appetitive Targeting scenario (not shown), all four variants perform statistically indistinguishably ($91.7$--$100\%$, $p=1.000$ for all three comparisons). (b) Attitude stability: Roll RMS increases significantly for all three ablations relative to the intact circuit (\emph{no lateral inhibition}: $+43.6\%$, $p<0.001$; \emph{no STN/STR}: $+45.9\%$, $p<0.001$; \emph{threshold only}: $+22.0\%$, $p=0.017$). Pitch RMS shows a different pattern: it increases for \emph{threshold only} ($+6.4\%$, $p=0.0029$), consistent with the roll-instability trend, but decreases for \emph{no lateral inhibition} and \emph{no STN/STR} ($-11.3\%$ and $-5.7\%$, both $p<0.005$) --- plausibly reflecting reduced net locomotor activity in trials that fail outright rather than improved attitude control.}
    \label{fig:ablation_combined}
\end{figure*}

Decision latency under conflict is unaffected by removing basal-ganglia mechanism ($81.7$~ms, 95\%~CI $[56.9, 106.4]$, for \emph{no lateral inhibition}, $p=0.34$; $66.6$~ms, 95\%~CI $[41.9, 91.2]$, for \emph{no STN/STR}, $p=0.84$; vs.\ $75.9$~ms, 95\%~CI $[59.3, 92.6]$, for \emph{full}), but is significantly slower under \emph{threshold only} ($109.1$~ms, 95\%~CI $[89.2, 129.0]$, $p=0.014$) --- consistent with a static threshold lacking the anticipatory, continuously-converging dynamics the WTA circuit provides during conflict resolution. Roll RMS increases significantly for all three ablations relative to the intact circuit ($+43.6\%$ \emph{no lateral inhibition}, $p<0.001$; $+45.9\%$ \emph{no STN/STR}, $p<0.001$; $+22.0\%$ \emph{threshold only}, $p=0.017$), indicating that attitude destabilization under conflict is a shared consequence of weakening the arbitration mechanism, regardless of which component is removed. Pitch RMS instead diverges by ablation type (Figure~\ref{fig:ablation_complex_stability}): it increases for \emph{threshold only} ($+6.4\%$, $p=0.0029$), consistent with the roll-instability trend, but decreases for \emph{no lateral inhibition} and \emph{no STN/STR} ($-11.3\%$ and $-5.7\%$, both $p<0.005$) --- plausibly a symptom of reduced net locomotor activity in the majority of trials that fail outright under these two ablations, rather than improved attitude control.

Taken together, these results support the winner-take-all arbitration design on two separable grounds: the basal-ganglia circuit's own mechanism, not merely the presence of \emph{some} priority rule, is what resolves genuine multi-stimulus conflict (R3-01); and no simpler substitute, including the same static-threshold family already used elsewhere in this architecture's Locomotion Module, matches it once continuous, real-time arbitration is what the task actually requires (R3-02).

\revblue{The study's sample size supports these two conclusions to different degrees. The first, that removing basal-ganglia mechanisms breaks conflict resolution, is robust to it; the second, that static threshold logic does not match the full circuit, is not yet conclusive. Each combination of variant and scenario comprises $n=12$ simulated trials, enough to detect large effects but with limited statistical power for smaller ones. Against the intact circuit's $12/12$ successes, Fisher's exact test at this sample size reaches $p<0.05$ only when an ablated variant fails in at least five of its twelve trials. The two mechanism-removing ablations fail in eight (\emph{no STN/STR}) and nine (\emph{no lateral inhibition}) of twelve trials, well beyond that boundary, so their collapse is not an artifact of the limited sample. The \emph{threshold only} variant's $83.3\%$ success rate ($10/12$, exact 95\%~CI $[51.6, 97.9]\%$), by contrast, is statistically inconclusive rather than evidence of equivalence with the intact circuit, and its decision-latency ($p=0.014$) and Roll RMS ($p=0.017$) increases, which lie closest to the significance boundary, should be regarded as provisional until replicated. Larger-scale repeated trials are left as future work to tighten these estimates and confirm the smaller effects. These would include more replicates per condition, a wider range of conflict configurations and noise levels, and an equivalent ablation protocol on the physical platform, which the present study does not include.}

\subsubsection{Response to Variable Topographies}

For the experiments involving variable and unstable topographies in simulation, a three‐section terrain map was designed, as shown in Figure~\ref{fig:TERR_IMU}. In the first simulated experiment, the robot begins moving forward driven by the hunting instinct encoded in the locomotor decision network. Subsequently, the transition from regular to rough terrain is introduced to analyze the standard deviation of the acceleration along the $Z$ axis throughout the simulation. \revblue{In the physical robot, this information corresponds to the MLP outputs $N_0$--$N_2$ (Section~\ref{ssec:Gait_Decision_Module}).} 

In a second test, the transition from flat terrain to an inclined slope is defined to evaluate the platform’s response to a sudden change in ground inclination. In this configuration, both in Gazebo and in the real robot, an aversive stimulus is positioned behind the platform; according to the previously defined behaviors, the robot must avoid it by switching to the quadrupedal mode. However, due to the slope, the platform adopts the differential rolling mode, achieving a more stable and efficient displacement. Note that while the simulated sloped terrain visually approximates a continuous curved trajectory in the Gazebo environment, the digital model (STL mesh) discretizes continuous geometry into planar facets. Consequently, the IMU telemetry in simulation captures these piecewise linear segments during motion, manifesting as alternating phases of dynamic angular change and transient stationary plateaus in the pitch response plots.

In the first simulation experiment shown in Figure~\ref{fig:RES_IMU}, the platform transitions immediately after entering the unstable terrain, adopting an appropriate gait to maintain efficient displacement. This behavior is reflected in the neural responses (Figure~\ref{fig:NEU_IMU}), where at 16 seconds, the unit $X_{0}$ shows the highest activity due to the large variation in vertical acceleration, a direct indicator of terrain irregularities. This phenomenon is also visible in Figure~\ref{fig:GRAPH_IMU}, where the standard deviation of the robot's vertical acceleration is plotted over time. The graph highlights the instant when this value exceeds the 3.7 threshold defined in the neural dynamics.

% RES_IMU merged with RES_IMU2 below, after the second-scenario prose (see fig:res_imu_pair).

\begin{figure}[t]
    \centering
    \includegraphics[width=0.5\textwidth]{NEU_IMU.png}
    \caption{Neural activations during locomotion over irregular terrain (simulation). (a)~Locomotion decision module; (b)~Gait mode control module. \revblue{Single trial ($N=1$).}}
    \label{fig:NEU_IMU}
\end{figure}

% GRAPH_IMU merged with GRAPH_IMU2 below, after the second-scenario prose (see fig:graph_imu_pair).

For the second stability test in simulation, the neural activities of the gait-mode control structure were recorded (Figure~\ref{fig:NEU_IMU2}b). \revblue{From $t \approx 38$~s, neuron $X_{1}$, which responds to pitch above the $U_{Pitch}$ threshold, excites unit $X_{16}$; $X_{16}$ suppresses $X_{15}$ and commands the switch from quadrupedal to differential rolling mode.} The transition during locomotion is illustrated in Figure~\ref{fig:RES_IMU2}, with snapshots of the platform posture taken every 12 seconds from a lateral perspective. It is important to highlight the neural responses shown in Figure~\ref{fig:NEU_IMU2}a, where the persistent activation of neuron $X_{4}$ produces sustained activity in neurons $X_{11}$ and $X_{12}$, which in turn drive neuron $X_{13}$ and ultimately generate the forward escape behavior. In Figure~\ref{fig:GRAPH_IMU2}, the data recorded by the inertial sensor are shown.

\begin{figure}[t]
    \centering
    \includegraphics[width=0.4\textwidth]{NEU_IMU2.png}
    \caption{Neural activations during locomotion over inclined terrain (simulation). (a)~Locomotion decision module; (b)~Gait mode control module. \revblue{Single trial ($N=1$).}}
    \label{fig:NEU_IMU2}
\end{figure}

\begin{figure}[t]
  \centering
  \begin{subfigure}[t]{\textwidth}
    \centering
    \includegraphics[width=0.85\linewidth]{RES_IMU.png}
    \caption{Trajectory in an unstable terrain in a simulation trial. The poses were taken every 22 seconds.}
    \label{fig:RES_IMU}
  \end{subfigure}

  \vspace{0.6em}

  \begin{subfigure}[t]{\textwidth}
    \centering
    \includegraphics[width=0.85\linewidth]{RES_IMU2.0.png}
    \caption{Trajectory during traversal of elevated terrain in simulation. Snapshots taken every 12~s.}
    \label{fig:RES_IMU2}
  \end{subfigure}
  \caption{Simulated-trajectory snapshot sequences for the two variable-topography scenarios. \revblue{Each panel shows a single trial ($N=1$ per panel).}}
  \label{fig:res_imu_pair}
\end{figure}

\begin{figure}[t]
  \centering
  \begin{subfigure}[t]{0.48\textwidth}
    \centering
    \includegraphics[width=\linewidth]{GRAPH_IMU.png}
    \caption{Standard deviation of the acceleration in the Z axis during transport on an unstable terrain.}
    \label{fig:GRAPH_IMU}
  \end{subfigure}\hfill
  \begin{subfigure}[t]{0.48\textwidth}
    \centering
    \includegraphics[width=\linewidth]{GRAPH_IMU2.png}
    \caption{Pitch angle during motion in a steeply sloping terrain.}
    \label{fig:GRAPH_IMU2}
  \end{subfigure}
  \caption{IMU signal traces motivating the mode transition in each simulated terrain scenario: (a) unstable/irregular terrain; (b) steeply sloping terrain. \revblue{Each panel shows a single trial ($N=1$ per panel).}}
  \label{fig:graph_imu_pair}
\end{figure}

\subsubsection{Quantitative Performance in Variable Topographies}
\label{sssec:variable_topo_quant}

\begin{figure*}[t]
    \centering
    \begin{subfigure}[t]{0.32\textwidth}
        \centering
        \includegraphics[width=\textwidth]{FigA_Stability_Profile.png}
        \caption{Dynamic stability profile (Tipover Risk, $TR$) aligned around the neural mode-switch instant ($t = 0.0\,\text{s}$) for  Rugged Terrain, upper panel and Inclined Slope, lower panel. Solid lines represent the mean across 15 independent trials; shaded bands denote $\pm 1\sigma$. The horizontal dashed line at $TR = 1.0$ marks the physical overturning boundary. No trial violated this limit across all 30 runs.}
        \label{fig:stability_profile}
    \end{subfigure}
    \hfill
    \begin{subfigure}[t]{0.32\textwidth}
        \centering
        \includegraphics[width=\textwidth]{FigB_Decision_Delay.png}
        \caption{Empirical Cumulative Distribution Functions (ECDF) of mechanical transition time ($T_{\text{switch}}$) across 15 randomized trials per terrain: (A)~Rugged Terrain, (B)~Inclined Slope. The mechanical hand-off is sub-millisecond in the majority of trials but shows a long tail to $\approx 0.6\,\text{s}$ in 2/15 rugged-terrain trials (Panel A) and a five-to-six-step spread across $0.9$--$1.8\,\text{ms}$ on the inclined slope (Panel B). Neural response latency ($T_{\text{response}}$) is not plotted here, as it was exactly $0.0\,\text{s}$ for all 30 trials with no dispersion to display — see the accompanying prose for that result.}
        \label{fig:decision_delay}
    \end{subfigure}
    \hfill
    \begin{subfigure}[t]{0.32\textwidth}
        \centering
        \includegraphics[width=\textwidth]{FigC_Terrain_Adaptability.png}
        \caption{Kinematic phase-space scatter plot of chassis Roll RMS vs. Pitch RMS for all 30 trials. Blue squares (Rugged Terrain) and orange triangles (Inclined Slope) separate clearly along the Pitch RMS axis, with overlapping Roll RMS ranges, demonstrating that the neural controller synthesizes a distinct, terrain-specialized locomotive signature for each environment. The dashed horizontal line marks the Pitch RMS midpoint ($\approx 1.98^{\circ}$) between the two terrain means, dividing the plot into the labeled Rugged Terrain and Inclined Slope zones — deliberately not drawn on the Roll RMS axis, where the overlap between terrains is itself the evidence of comparably effective lateral stabilization (see accompanying prose).}
        \label{fig:terrain_adaptability}
    \end{subfigure}
    \caption{Quantitative performance profiles for the variable-topography experimental battery (15 trials each terrain type). (A)~Dynamic Stability and Tipover Risk profile aligned at the mode-switch instant. (B)~Decision Delay via ECDF analysis of mechanical transition time (neural response latency is reported numerically in the text, not plotted — see Section~\ref{sssec:variable_topo_quant}). (C)~Terrain Adaptability Phase Space mapping chassis Roll RMS against Pitch RMS. Together, these three profiles provide the quantitative evidence base for evaluating the bio-inspired controller's safety, speed, and kinematic adaptability under stochastic variable-topography conditions.}
    \label{fig:variable_topo_metrics}
\end{figure*}

\revblue{The Tipover Risk ($TR$) index was tracked across 15 randomized runs per terrain (Figure~\ref{fig:stability_profile}).} The $TR$ metric was formally derived from the McGhee--Frank Normalized Stability Margin as $TR = 1 - SM/r_{in}$, where $SM$ is the shortest horizontal distance between the center of mass projection and the support polygon boundary, and $r_{in}$ is the inradius of that polygon. Values approaching $TR = 0.0$ represent maximum structural safety, while $TR = 1.0$ marks the critical overturning threshold. In both terrains, the platform maintained a consistently bounded risk profile around $\mu_{TR} \approx 0.58$--$0.60$ (rugged terrain mean $0.600$; inclined slope mean $0.593$) throughout the recorded window. Rather than a monotonic escalation, the rugged-terrain trace showed an isolated transient spike reaching $TR \approx 0.77$ shortly before the mode-switch instant ($t \approx -1.3$ to $-0.5\,\text{s}$), consistent with the periodic contraction of the support polygon as individual legs were lifted during the quadrupedal gait cycle, before settling back to the $\approx 0.59$ baseline across the transition at $t = 0.0\,\text{s}$. The inclined-slope trace remained flat at $TR \approx 0.59$ across the entire window, with no escalation or drop coincident with the mode switch. \revblue{Across all 30 trials the $\pm 1\sigma$ envelope never crossed $TR = 1.0$.}


To quantify the computational efficiency and temporal determinism of the proposed bio-inspired control architecture under stochastic conditions, the neural response latency ($T_{\text{response}}$) and the Empirical Cumulative Distribution Function (ECDF) of the mechanical state-switching delay ($T_{\text{switch}}$) were analyzed across 15 randomized trials for both experimental terrains, the latter illustrated in Figure~\ref{fig:decision_delay}. \revblue{In these trials, zero-mean Gaussian noise proportional to the measured value was added to the IMU, LiDAR, and camera inputs at five levels ($\sigma = 0$, 5, 10, 20, and 30\%), with three trials per level on each terrain.} \revblue{The neural latency was} $T_{\text{response}} = 0.0\,\text{s}$ for both terrain types, showing that the threshold-based basal ganglia selection logic filtered out the injected sensor noise and dispatched the correct locomotion command within the same simulation frame as the environmental disturbance. The mechanical transition latency was likewise sub-millisecond in the vast majority of trials, but not fully invariant across terrain types. On rugged terrain, $T_{\text{switch}}$ fell within $0.2$--$0.7\,\text{ms}$ for 13 of 15 trials (a single discrete update cycle of the ROS~2 Control loop operating at $3.3\,\text{kHz}$), while the remaining 2 trials ($13.3\%$) exhibited a long tail out to $\approx 0.6\,\text{s}$, indicating that occasional trials require additional update cycles to complete the locomotion-matrix hot-swap under this terrain's contact conditions. On the inclined slope, $T_{\text{switch}}$ spread across $0.9$--$1.8\,\text{ms}$ in a five-to-six-step staircase rather than collapsing onto a single value, reflecting genuine trial-to-trial dispersion of a comparable order of magnitude rather than a perfectly deterministic delay. \revblue{Neural decision dispatch thus remained deterministic across all 30 trials, while the residual mechanical variability, including the rugged-terrain long tail, is a bounded, terrain-dependent effect rather than chattering.}

\revblue{The steady-state kinematics of both terrains were mapped onto a Roll RMS ($\phi_{\text{RMS}}$) vs.\ Pitch RMS ($\theta_{\text{RMS}}$) phase space across all 30 successful trials (Figure~\ref{fig:terrain_adaptability}).} The results revealed a clear separation in Pitch RMS between the two terrains, with partial overlap in Roll RMS. The rugged terrain converged into a low-to-moderate roll, low-pitch operational zone ($\mu_{\phi} \approx 1.59^{\circ}$, $\mu_{\theta} \approx 1.67^{\circ}$), capturing the kinematic footprint of asymmetric foot-ground interactions on uneven stones while maintaining a largely level chassis. Conversely, the inclined slope clustered in a higher-pitch region with comparable roll magnitude ($\mu_{\phi} \approx 1.47^{\circ}$, $\mu_{\theta} \approx 2.29^{\circ}$). This spatial distribution reflects the geometric adaptation triggered by the Hybrid morphological transformation: during slope ascent, the chassis aligned longitudinally with the ramp gradient (higher Pitch RMS), while the deployment of the coaxial wheel base kept the lateral support polygon wide enough that roll disturbances remained similar in magnitude to the rugged-terrain baseline rather than being further suppressed. \revblue{The Pitch-RMS separation indicates that the network selects a distinct, repeatable posture for each terrain rather than applying an environment-agnostic control law.}

\subsubsection{Quantitative Stability Analysis of Locomotion Mode Transitions}
\label{sssec:transition_stability}

%
\revblue{This subsubsection quantifies static stability for the Quadrupedal Walking gait and all six pairwise transitions among the three locomotion modes, using the Normalized Stability Margin ($SM$) and Tipover Risk ($TR$) of Section~\ref{sssec:variable_topo_quant}, computed frame-by-frame from the foot-contact positions logged in simulation.}


\paragraph{Quadrupedal Walking gait stability.}

%
The Quadrupedal Walking gait is intrinsically
static-stable by construction: at every instant, at least three feet
maintain ground contact, guaranteeing a non-degenerate
support polygon throughout the full cycle.
Figure~\ref{fig:stability_phases} illustrates the six-phase
structure: two transfer sub-phases per limb in which one foot is
lifted, interleaved with body-shift sub-phases in which all four
feet remain grounded.
Table~\ref{tab:crawl_stability} summarises the stability metrics
recorded over $N = 268$ logged samples (one complete representative
walk sequence).


\begin{figure}[t]
    \centering
    \includegraphics[width=0.85\textwidth]{fig_stability_phases.png}
    \caption{%
        %
        Six-phase gait cycle of the Quadrupedal Walking mode ($C$).
        Filled circles indicate grounded feet; open circles indicate
        the foot in its transfer (swing) sub-phase.
        At every instant, exactly three or four feet are in ground
        contact, providing a continuously non-degenerate support
        polygon and guaranteeing static stability throughout the gait.
        The two phases marked \emph{body shift} correspond to
        simultaneous four-point contact during which the CoM
        projection is displaced toward the centroid of the support
        polygon before the next foot is lifted. \revblue{Data: $N=268$ logged samples from one complete representative walk sequence (Table~\ref{tab:crawl_stability}).}%
        %
    }
    \label{fig:stability_phases}
\end{figure}

\begin{table}[t]
\centering
\caption{%
    Normalized Stability Margin ($SM$), Tipover Risk ($TR$), and
    support-polygon inradius during Quadrupedal Walking
    ($N = 268$ logged samples).
    Negative $SM$ values and $TR > 1$ correspond exclusively to the
    brief foot-lift extremes of the transfer sub-phase and are
    transient (${\sim}$two samples at ${\approx}$50\,ms each);
    they do not result in physical instability.%
}
\label{tab:crawl_stability}
\begin{tabular}{lcccc}
\toprule
\textbf{Metric} & \textbf{Mean} & \textbf{Std Dev} &
\textbf{Min} & \textbf{Max} \\
\midrule
$SM$ (m)       & 0.034 & 0.018 & $-$0.015 & 0.050 \\
$SM_{\mathrm{norm}}$ & 0.561 & 0.409 & $-$0.173 & 1.000 \\
$TR$           & 0.439 & 0.409 & 0.000   & 1.173 \\
Inradius (m)   & 0.054 & 0.049 & $\approx 0$ & 0.101 \\
\bottomrule
\end{tabular}
\end{table}

%
The mean normalized margin ($SM_{\mathrm{norm}} = 0.561 \pm 0.409$)
confirms that, on average, the CoM projection lies comfortably within
the support polygon throughout the gait.
The transient excursions above $TR = 1$ are an artefact of the
three-point polygon contracting to its minimum area at the instant
the lifted foot reaches its farthest swing position; the low inertia
of the slow gait cadence ensures that these short-duration events
do not result in actual tip-over events.


\paragraph{Mode-transition stability.}

%
All six pairwise transitions among the three locomotion
modes ---Quadrupedal Walking ($C$), Differential Rolling ($H$),
and Omnidirectional ($X$)--- were executed sequentially in a
single continuous simulation run and recorded with full
foot-contact and CoM telemetry ($N_{\mathrm{total}} = 979$ samples
across the complete sequence).
\revblue{Both directions of each mode pair converge to statistically indistinguishable support-polygon geometries and $TR$ values, so Figure~\ref{fig:transition_phases} shows each undirected pair ($C \leftrightarrow H$, $H \leftrightarrow X$, $X \leftrightarrow C$) once.}

\begin{figure}[t]
    \centering
    \includegraphics[width=\textwidth]{fig_transition_phases_v10.png}
    \caption{%
        Static stability across reversible locomotion mode transitions,
        one row per undirected mode pair. For each pair, the two columns
        show the settled support-polygon geometry in each of the two
        modes; the double-headed arrow denotes that the transition is
        equally stable in both directions. $TR$ and $SM$ are reported as
        a min--max range across all instances of that mode observed
        across the original six directed transitions (e.g. mode $C$'s
        geometry recurs as the onset of $C\to H$, the final pose of
        $X\to C$, the onset of $C\to X$, and the final pose of
        $H\to C$) rather than as a single value per direction, since all
        four readings of a given mode agree closely --- itself the
        empirical evidence that both transition directions are equally
        stable. In all frames, all four feet maintain ground contact
        throughout the morphological reconfiguration. \revblue{Data: a single continuous run covering all six directed transitions ($N_{\mathrm{total}}=979$ samples).}%
    }
    \label{fig:transition_phases}
\end{figure}

%
A critical property of the mode-switching protocol is that
morphological reconfigurations are constrained to states in which
all four feet simultaneously maintain ground contact.
Crucially, this constraint does not require the robot to halt:
the transition command may be issued while the chassis is actively
translating or rotating, and locomotion in the target mode resumes
immediately upon completion of the limb reconfiguration, without an
intermediate stop phase.
Across all six pairwise transitions ---$C \to H$, $H \to X$,
$X \to C$, $C \to X$, $X \to H$, and $H \to C$--- and all
logged samples, every transition completed
successfully, yielding a transition success rate of $6/6$ ($100\%$).
At the instant of morphological switching itself --- the
moment all four feet share ground contact for maximum support --- the
Normalized Stability Margin reached its platform-geometry maximum
($SM = 0.050\,\text{m}$, $SM_{\mathrm{norm}} = 1.000$) and the
Tipover Risk index reached $TR = 0.000$, without exception; this
instantaneous reading does not describe the full $N_{\mathrm{total}} =
979$-sample window, which also includes established-gait locomotion
between switching events, where $TR$ takes non-zero values (see
Figure~\ref{fig:transition_phases}'s reported per-mode ranges).


%
\revblue{Because reconfiguration is performed with all four feet in contact, no abrupt postural discontinuity is introduced at the switching instant.}
The steady-state Roll and Pitch RMS bounds reported in
Figures~\ref{fig:variable_topo_metrics}
and Table~\ref{tab:consolidated_simulation_metrics}
characterize the locomotion phases immediately preceding and
following each transition; \revblue{the mean Roll RMS stays at or
below $1.59^\circ$ and the mean Pitch RMS at or below $2.29^\circ$
across these suites and both variable-topography terrains},
indicating that the mode-switching protocol does not introduce
roll or pitch instabilities in either the source or target
locomotion phase.


%
\revblue{Joint-level energy was not logged in Gazebo (no power-monitoring plugin); energy per mode is characterized on the physical platform in Section~\ref{ssec:QuantReal}.}


\subsubsection{FSM Gait Arbitration Baseline Comparison}
\label{sssec:fsm_baseline_comparison}

\revblue{Table~\ref{tab:fsm_comparison} compares the FSM controller of Section~\ref{ssec:FSMBaseline} with the basal-ganglia network on all seven simulated suites, reporting success rate, task completion time, Roll RMS, and $\lambda$, a motor-command clarity ratio (commanded-velocity magnitude relative to neural activity norm; higher indicates more decisive output).}

\begin{table}[htbp]
\centering
\caption{FSM vs.\ basal-ganglia (Neural) gait arbitration, all seven simulated suites. Complex/FSM has no completed trials (0\% success), so time/Roll RMS are undefined ($-$). One Obstacle/Neural trial with an anomalous roll reading (physically implausible per-timestep values, not a measurement of stable attitude) was excluded from the Roll RMS/task-time statistics below, consistent with Table~\ref{tab:consolidated_simulation_metrics}; its own success/failure verdict is unaffected. \revblue{$N$ is the number of independent missions behind each success rate (Appetitive, Aversive, Obstacle, Complex; a mission that first times out is retried until it completes) or of individual attempts (Inspection, Rugged Terrain, Inclined Slope); for Rugged Terrain and Inclined Slope the Neural controller was evaluated only without perturbation, hence its small $N$.}}
\label{tab:fsm_comparison}
\begin{tabularx}{\textwidth}{l l >{\centering\arraybackslash}X >{\centering\arraybackslash}X >{\centering\arraybackslash}X >{\centering\arraybackslash}X >{\centering\arraybackslash}X}
\toprule
\textbf{Family} & \textbf{System} & \revblue{\textbf{N}} & \textbf{SR (\%)} & \textbf{Task Time (s)} & \textbf{Roll RMS (deg)} & \textbf{$\lambda$} \\
\midrule
Appetitive & Neural & \revblue{15} & 100.0  & 38.52 $\pm$ 1.20  & 0.536 $\pm$ 0.009 & 0.990 \\
Appetitive & FSM    & \revblue{15} & 100.0 & 38.96 $\pm$ 1.77  & 0.525 $\pm$ 0.003 & 0.952 \\
Aversive   & Neural & \revblue{15} & 100.0  & 9.98 $\pm$ 2.07   & 1.357 $\pm$ 0.106 & 0.925 \\
Aversive   & FSM    & \revblue{15} & 100.0 & 8.32 $\pm$ 1.08   & 1.825 $\pm$ 0.112 & 0.731 \\
Obstacle   & Neural & \revblue{15} & 100.0 & 15.72 $\pm$ 2.49  & 1.547 $\pm$ 0.154 & 0.945 \\
Obstacle   & FSM    & \revblue{15} & 100.0 & 12.44 $\pm$ 1.09  & 1.499 $\pm$ 0.089 & 0.811 \\
Complex    & Neural & \revblue{15} & 100.0 & 24.96 $\pm$ 10.36 & 1.034 $\pm$ 0.349 & 0.976 \\
Complex    & FSM    & \revblue{15} & 0.0   & --                & --                & 0.942 \\
Inspection & Neural & \revblue{20} & 85.0  & 51.61 $\pm$ 11.96 & 1.302 $\pm$ 0.078 & 1.000 \\
Inspection & FSM    & \revblue{20} & 45.0  & 55.13 $\pm$ 14.25 & 1.270 $\pm$ 0.136 & 0.985 \\
Rugged Terrain & Neural & \revblue{6} & 50.0  & 54.94 $\pm$ 2.72  & 1.827 $\pm$ 0.033 & 0.909 \\
Rugged Terrain & FSM    & \revblue{25} & 60.0  & 55.42 $\pm$ 8.56  & 1.739 $\pm$ 0.140 & 0.891 \\
Inclined Slope & Neural & \revblue{3} & 100.0 & 77.02 $\pm$ 1.26  & 6.751 $\pm$ 7.160 & 0.996 \\
Inclined Slope & FSM    & \revblue{21} & 71.4  & 64.53 $\pm$ 2.79  & 5.430 $\pm$ 8.209 & 0.993 \\
\bottomrule
\end{tabularx}
\end{table}

In the three single-stimulus suites (Appetitive, Aversive, Obstacle), the FSM's fixed priority rule is structurally well-matched to the task --- one stimulus is active at a time, so a static hierarchy resolves it immediately and deterministically. Both controllers eventually complete every mission in all three suites (100\% mission-level success in Table~\ref{tab:fsm_comparison}, i.e.\ a trial that first times out and is then retried still counts as a success once it lands, and this is the only sense in which "success rate" is used anywhere in this paper). \revblue{The two controllers diverge, however, in how often that retry is needed (Figure~\ref{fig:attempts_vs_noise_level}; attempts needed per completed mission, not a success rate):} Neural needs up to $1.67\times$ as many attempts on average in Appetitive and $1.33\times$ in Aversive at the noisiest levels tested, non-monotonically given the small per-level sample, while FSM stays at exactly one attempt throughout.

\begin{figure}[t]
    \centering
    \includegraphics[width=\textwidth]{fig_attempts_vs_noise_level.png}
    \caption{Attempts needed per completed mission vs.\ noise level, single-stimulus suites (not a success-rate metric --- every mission in this dataset eventually succeeds, see Table~\ref{tab:fsm_comparison}; this counts how many attempts, including retries after a timeout, that took). FSM never needs a retry at any level in any of the three panels (flat at $1.0$). Neural needs an occasional retry at intermediate/high noise in Appetitive and Aversive, non-monotonically (each point averages only 3 trials --- read point-to-point zig-zags as sampling variance, not a mechanistic trend). In the Obstacle panel both systems sit at $1.0$ across all five evaluated noise levels --- not a subset --- so the dashed Neural line sits exactly on top of the solid FSM line.}
    \label{fig:attempts_vs_noise_level}
\end{figure}

This advantage reverses in the Complex suite, where three stimuli (obstacle, aversive, appetitive) must be handled in sequence: Neural reaches 100\% success while FSM reaches 0\% (all 15 trials timing out or, in one case, tipping over). The mechanism is not that the FSM decides worse while a conflict is active, but that it stays in an unresolved simultaneous conflict (both aversive and appetitive stimuli detected at once) far longer before its fixed priority rule forces a choice: real-time conflict exposure averages $\approx$0.2~s/trial (1.7\% of trial time) for Neural against $\approx$6.1~s/trial (4.7\%) for FSM --- roughly 30$\times$ longer in absolute time. Once the FSM's obstacle-avoidance rule releases control, the robot is frequently repositioned within a second stimulus's detection range, and in the absence of any active rule the controller defaults to an undirected exploratory mode with no residual bias toward the original target, unlike the Neural system's persistent-activation memory (Section~\ref{ssec:FSMBaseline}). \revblue{Figure~\ref{fig:mode_timeline_complex} contrasts representative mode traces of both controllers in this suite.}

\begin{figure}[t]
    \centering
    \includegraphics[width=\textwidth]{fig_mode_timeline_complex.png}
    \caption{Active gait mode over time, Complex suite: a representative successful Neural trial (top) vs.\ a representative failed FSM trial (bottom, timeout). The FSM trace switches abruptly between modes and holds each one flat for tens of seconds at a time, since its fixed-duration persistence timers either hold a mode or release it outright with no intermediate state; the Neural trace shows shorter, more graduated occupancy of each mode, reflecting the continuous decay of the basal ganglia's lateral inhibition rather than a hard timer. \revblue{Each trace is the first trial by index (not hand-picked) of $N=15$ trials per controller.}}
    \label{fig:mode_timeline_complex}
\end{figure}

On the two terrain suites (Rugged Terrain, Inclined Slope), success rate, task time, and Roll RMS are reported in Table~\ref{tab:fsm_comparison}, but $T_{switch}$/$T_{response}$ are deliberately omitted: the team that ran the FSM terrain trials identified three configuration differences from the basal-ganglia code's terrain-transition logic (the pitch-instability threshold, a transition timing offset, and the switching-time formula's margin term) that make those two specific metrics not comparable under strictly identical conditions without re-running both controllers with matched parameters, which was out of scope for this comparison. The Neural side of this comparison also draws on a substantially smaller sample (six Rugged Terrain trials, three Inclined Slope trials, restricted to the zero-noise condition to match the FSM runs) than the FSM side (25 and 21 trials respectively); the apparent tipover-rate difference on the Inclined Slope (zero tipovers in the small Neural sample vs.\ six of 21 FSM trials) is reported as an observation rather than a statistically supported finding given this sample-size imbalance.

\revblue{The Neural-side values for the Appetitive, Aversive, Obstacle, and Complex suites are recomputed from the raw trial logs and match Table~\ref{tab:consolidated_simulation_metrics} on every reported statistic, under the single success-rate definition given above.}

\FloatBarrier

\subsubsection{Reactive Inspection Task}
\label{sssec:reactive_inspection}

\revblue{To obtain inspection-oriented task indicators, the platform was evaluated in the dedicated obstacle-field scenario described in Section~\ref{ssec:Gazebo_Simulation}, in which the target is not visible from the initial pose (sampling rule in Table~\ref{tab:inspection_metrics}).}

\revblue{As described in Section~\ref{ssec:Gazebo_Simulation}, behavior in this task is purely reactive (no map, no spatial memory); re-routing emerges from the basal-ganglia/winner-take-all (BG/WTA) arbitration network integrating LiDAR and camera input in real time.} Coverage rate --- the percentage of a predefined inspection area visited --- is not reported here, as it presupposes a coverage planner and a known target area that this reactive architecture neither has nor requires. Instead, task-completion indicators that are both measurable with the existing instrumentation and consistent with the system's actual capabilities are used: success rate (SR), final approach precision ($d_{final}$), and time to first visual acquisition of the target ($T_{acquisition}$).

$d_{final}$ is the Euclidean distance between the robot and the target at the instant the inspection-stop neuron activates ($X_{17} > 0.25$, sustained 2~s); $N_{ev}^{LiDAR}$ counts sensory-level obstacle-detection transitions (LiDAR channel crossing its 0.3 activation threshold); $N_{ev}^{X}$ counts the resulting omnidirectional-mode switches, gated by a 3~s dwell-time hysteresis; $T_{acquisition}$ is the simulation time from post-warmup to the target's bounding-box area first exceeding the visibility threshold.

Twenty trials yielded SR $= 85\%$ (17/20), with three timeouts (trials 7, 11, 17) discussed below. Among the 17 successful trials, $d_{final} = 1.076 \pm 0.038$~m (Table~\ref{tab:inspection_metrics}): the BG$\rightarrow X_{17}$ loop converges to essentially the same stopping distance from the target regardless of the route taken to reach it.

\begin{table}[htbp]
\centering
\caption{Reactive Inspection Task --- Aggregate Performance Metrics (17/20 successful trials, obstacle field with initially non-visible target). \revblue{Simulation data (no physical battery). Sampling rule: 20 independent trials, 90\,s timeout, no noise injection; mean $\pm$ sample SD, min and max over the successful trials only.}}
\label{tab:inspection_metrics}
\begin{tabularx}{\textwidth}{l >{\centering\arraybackslash}X >{\centering\arraybackslash}X >{\centering\arraybackslash}X}
\toprule
\textbf{Metric} & \textbf{Mean $\pm$ SD} & \textbf{Min} & \textbf{Max} \\
\midrule
$d_{final}$ (m)              & 1.076 $\pm$ 0.038 & 0.991 & 1.130 \\
$N_{ev}^{LiDAR}$          & 7.71 $\pm$ 6.42   & 4     & 28    \\
$N_{ev}^{X}$        & 3.29 $\pm$ 2.80   & 1     & 10    \\
$T_{acquisition}$ (s)        & 32.2 $\pm$ 14.1   & 20.7  & 70.8  \\
$T_{task}$ (s)               & 51.6 $\pm$ 12.3   & 42.8  & 86.8  \\
\bottomrule
\end{tabularx}
\\[2pt]
{\footnotesize Success rate: 17/20 (85\%); three timeouts at 90~s (trials 7, 11, 17), analyzed below.}
\end{table}

The 17 successful trials separate into two navigation patterns distinguished by $N_{ev}^{LiDAR}$: a clean pattern ($n=12$, $N_{ev}^{LiDAR} = 4$--$7$, close to one detection per capsule, $d_{final} = 1.081 \pm 0.042$~m) and an elevated-evasion pattern ($n=5$, $N_{ev}^{LiDAR} = 11$--$28$, additional evasive maneuvering before finding a route toward the target, $d_{final} = 1.065 \pm 0.022$~m). Despite the substantially messier routes, the elevated-evasion group converges at least as tightly on the stopping point as the clean group (Figure~\ref{fig:inspection_dfinal_scatter}) --- if anything, more tightly, though the small elevated-group sample ($n=5$) cautions against over-reading that difference. The three timeouts are extreme cases of the elevated pattern ($N_{ev}^{LiDAR} = 18$, $27$, $33$): in all three, the target was visually detected ($T_{acquisition}$ non-null) but the appetitive channel never accumulated sufficient sustained activation to trigger $X_{17}$ before the deadline, as competing obstacle-channel activation persisted throughout. This is consistent with the architecture's known limitation --- reactive navigation without spatial memory can fail under unfavorable dynamic configurations --- rather than a failure of the inspection-stop mechanism itself.

\begin{figure}[t]
    \centering
    \includegraphics[width=0.7\textwidth]{fig_inspection_dfinal_vs_nlidar.png}
    \caption{Final approach precision ($d_{final}$) versus sensory evasion-event count ($N_{ev}^{LiDAR}$) across the 17 successful Reactive Inspection Task trials, split into a clean-traversal pattern ($N_{ev}^{LiDAR} < 10$, $n=12$, $d_{final} = 1.081 \pm 0.042$~m) and an elevated-evasion pattern ($N_{ev}^{LiDAR} \geq 10$, $n=5$, $d_{final} = 1.065 \pm 0.022$~m). \revblue{Simulation data; same sampling rule as Table~\ref{tab:inspection_metrics}.}}
    \label{fig:inspection_dfinal_scatter}
\end{figure}

\revblue{This experiment also quantifies re-routing after obstacle blockage:} \revblue{behavioral mode changes ($N_{ev}^{X}$) were $\approx 2.3\times$ fewer than the sensory-level obstacle detections ($N_{ev}^{LiDAR}$) that triggered them (Table~\ref{tab:inspection_metrics})}, in every one of the 17 successful trials without exception. This demonstrates that the WTA arbitration converts repeated raw sensory fluctuations into a smaller number of sustained behavioral commitments, avoiding the mode-dithering a simple threshold-triggered controller would exhibit --- the mechanism underlying the reactive re-routing observed in this task.

\revblue{$T_{acquisition}$ (Table~\ref{tab:inspection_metrics})} is a first-acquisition metric, not re-acquisition, since the target is not visible from the initial pose in this scenario. Visual pursuit is itself conditioned on the absence of a frontal obstacle: during evasion it is suspended, resuming automatically once the obstacle clears --- the system's re-acquisition-like behavior. The robot remained in the wheeled/differential locomotion mode (H) throughout all 20 trials, so no gait-transition stability data applies to this experiment.

Under the FSM baseline (Section~\ref{ssec:FSMBaseline}) on this same scenario, success rate drops to 45\% (9/20) against the Neural system's 85\% (Table~\ref{tab:fsm_comparison}). The obstacle-avoidance rule's fixed persistence timer re-triggers whenever the FSM re-detects an obstacle after leaving avoidance mode, and in 11 of 20 trials the robot remains oriented such that it keeps re-entering avoidance, occupying that mode 93--99\% of trial time and never visually detecting the target at all --- a qualitative coverage failure, distinct from the Neural system's occasional timing-based misses (every Neural trial detects the target at least once). Figure~\ref{fig:mode_timeline_inspection} contrasts the resulting mode traces; the FSM's abrupt, timer-driven switching follows the same pattern described for the Complex suite (Section~\ref{sssec:fsm_baseline_comparison}).

\begin{figure}[t]
    \centering
    \includegraphics[width=\textwidth]{fig_mode_timeline_inspection.png}
    \caption{Active gait mode over time, Reactive Inspection Task: a representative successful Neural trial (top) vs.\ a representative successful FSM trial (bottom) --- shown here for a case where the FSM does escape avoidance mode; in its 11 failing trials the lower trace instead stays flat in avoidance for nearly the entire duration. Same qualitative pattern as Figure~\ref{fig:mode_timeline_complex}, reproduced here for this suite. \revblue{Each trace is the first trial by index (not hand-picked) of $N=20$ trials per controller.}}
    \label{fig:mode_timeline_inspection}
\end{figure}

\FloatBarrier

\subsection{Physical Implementation Results}

Complementarily, the experiments with the physical platform reproduced the same conditions as those used in simulation. Stability tests were conducted on inclined and irregular terrains, as described in subsection~\ref{ssec:Construction}. 

\subsubsection{Single Stimulus Response}
The same protocol was replicated on the physical robot; Figure~\ref{fig:RESPUESTAS_EST_UNI_real} summarizes the resulting trajectories for the three stimuli --- approach (a, appetitive), evasion (b, obstacle), and escape (c, aversive) --- all consistent with the simulation results.

\begin{figure}[t]
    \centering
    \includegraphics[width=0.8\linewidth]{RES_EST_UNI_real.png}
    \caption{Response to single stimulus in the real robot. \revblue{Each panel shows a single trial ($N=1$ per panel).}}
    \label{fig:RESPUESTAS_EST_UNI_real}
\end{figure}


In the obstacle-evasion response (Figure~\ref{fig:NEU_EST_UNIG_real}a), input-layer neurons $X_{15}$/$X_{16}$ activate first, indicating the obstacle on the front-right side; $L_{4}$ reaches maximum activity after $\approx35$~s, and Figure~\ref{fig:NEU_EST_UNIG_real}c shows unit $X_{14}$ engaging omnidirectional gait for the leftward evasive maneuver.

\begin{figure}[t]
    \centering
    \includegraphics[width=1.0\textwidth]{NEU_EST_UNIG_real.png}
    \caption{Raster plot of neural activations for single obstacle stimulus in real robot. \revblue{Single trial ($N=1$).}}
    \label{fig:NEU_EST_UNIG_real}
\end{figure}

During the 5~s locomotion-mode change (differential drive $\rightarrow$ omnidirectional), $X_{15}$/$X_{16}$ remain nearly constant (no displacement); after $\approx40$~s their activity decreases and shifts to neighboring neurons, marking successful completion of the maneuver.

For the appetitive stimulus, the physical robot exhibits a frontal-approach behavior: at $\approx23$~s, neurons $X_5$/$X_7$ (Figure~\ref{fig:NEU_EST_UNIG_real}b) drive a leftward orientation correction via $X_{11}$ to center the target. This differs from simulation (Figure~\ref{fig:NEU_EST_UNIB}), where $X_{12}$'s peak activity at the trial's start reflects an immediate rightward rotation, since the appetitive stimulus was laterally positioned in simulation but directly in front of the robot in the physical trial.


\subsubsection{Response to Multiple Stimuli of Different Natures}

In the physical experiment (Figure~\ref{fig:robot-6frames}), the input layer (Figure~\ref{fig:NEU_EST_MUL_real}b) shows simultaneous left-side activation from a LiDAR-detected object; since it is farther than the right-side obstacle, response/auxiliary-layer activity decreases progressively and the network disregards it, with unit $L_2$ only mildly activated, attenuated by $L_3$ inhibition.

\begin{figure}[t]
    \centering
    \includegraphics[width=1.0\textwidth]{NEU_EST_MUL_real.png}
    \caption{Raster plot of neural activations for multiple stimuli of diverse nature in real robot. \revblue{Single trial ($N=1$).}}
    \label{fig:NEU_EST_MUL_real}
\end{figure}

A simultaneous $X_{15}$/$X_{16}$ activation (Figure~\ref{fig:NEU_EST_MUL_real}d), from the constant hunting-instinct stimulus, keeps $X_{16}$ active and uninhibited by $X_{15}$ --- the no-motor-action state shown in Figure~\ref{fig:frames-c}.

\begin{figure}[t]
  \centering
  \begin{subfigure}[t]{0.48\textwidth}
    \centering
    \includegraphics[width=\linewidth]{D1.png}
    \caption{Initial state: all three stimulus classes in view.}
    \label{fig:frames-a}
  \end{subfigure}\hfill
  \begin{subfigure}[t]{0.48\textwidth}
    \centering
    \includegraphics[width=\linewidth]{D3.png}
    \caption{Indecisive hold ($X_{16}$ active, uninhibited).}
    \label{fig:frames-c}
  \end{subfigure}

  \vspace{0.6em}

  \begin{subfigure}[t]{0.48\textwidth}
    \centering
    \includegraphics[width=\linewidth]{D5.png}
    \caption{Evasive dispatch, $\approx46$~s.}
    \label{fig:frames-e}
  \end{subfigure}\hfill
  \begin{subfigure}[t]{0.48\textwidth}
    \centering
    \includegraphics[width=\linewidth]{D6.png}
    \caption{Goal approach and stop, $\approx50$--$90$~s.}
    \label{fig:frames-f}
  \end{subfigure}
  \caption{Four representative frames from the real-robot experiment facing visual stimuli of different nature: (a)–(d) show the temporal evolution across the trial, detailed in the accompanying text.}
  \label{fig:robot-6frames}
\end{figure}

Following the aversive globus pallidus's peak activation at $\approx46$~s, the robot exits the indecisive state into quadrupedal mode, with rightward-directionality neurons $X_8$/$X_{12}$ reaching maximum amplitude (Figure~\ref{fig:NEU_EST_MUL_real}c); the resulting evasive maneuver reveals the appetitive stimulus after $\approx50$~s, triggering approach. A leftward correction centers the stimulus (neurons $X_7$/$X_{11}$), followed by a final rightward readjustment; the robot then holds a straight trajectory, stopping at $\approx90$~s when $X_{17}$ fully inhibits movement.

\subsubsection{Evaluation of the MLP-based terrain classifier}
\label{sssec:mlp_eval}

To complement the qualitative description of the terrain-classification module (Section~\ref{ssec:Gait_Decision_Module}), the MLP (150–80 hidden units, LeakyReLU with negative slope $0.01$, Adam optimizer) was evaluated on the IMU dataset ($n = 600$ samples, 135-dimensional input, binary labels: 0 = flat terrain, 1 = inclined terrain) using a stratified 80/20 train–test split and 5-fold stratified cross-validation.

Across the 5-fold cross-validation, the model achieved an accuracy of $0.767 \pm 0.028$, precision of $0.764 \pm 0.024$, recall of $0.779 \pm 0.048$, and an F1-score of $0.771 \pm 0.033$. On the held-out test set ($n = 120$\revblue{; confusion matrix in Figure~\ref{fig:mlp_eval_grid}a}), the model reached an accuracy of 0.800, precision of 0.747, recall of 0.918, F1-score of 0.824, and an ROC-AUC of 0.864, confirming that the network reliably discriminates between flat and inclined terrain, with a recall bias toward the inclined-terrain class that minimizes the safety-critical risk of missed inclines. \revblue{The ROC curve (Figure~\ref{fig:mlp_eval_grid}b) stays well above the chance diagonal across the operating range.}

Figure~\ref{fig:mlp_eval_grid}c illustrates the training convergence behavior, showing the evolution of the learning loss across epochs and confirming stable optimization of the selected configuration. To justify the choice of architecture, an ablation study compared the proposed 150–80 configuration against simpler (50, 100) and deeper (200, 100, 50) topologies. As shown in Figure~\ref{fig:mlp_eval_grid}d, the 150–80 configuration achieved the highest F1-score among all tested architectures ($\text{F1} = 0.771$), outperforming the shallower (50: $\text{F1} = 0.747$; 100: $\text{F1} = 0.766$) and deeper (200-100-50: $\text{F1} = 0.766$) variants. This confirms that 150–80 represents the most favorable trade-off between model capacity and generalization for this classification task, while remaining computationally inexpensive for onboard deployment.

\begin{figure*}[t]
    \centering
    \begin{subfigure}[b]{0.48\textwidth}
        \centering
        \includegraphics[width=\linewidth]{fig1_confusion_matrix.png}
        \caption{Confusion matrix (test set).}
    \end{subfigure}
    \hfill
    \begin{subfigure}[b]{0.48\textwidth}
        \centering
        \includegraphics[width=\linewidth]{fig2_roc_curve.png}
        \caption{ROC curve (AUC = 0.864).}
    \end{subfigure}

    \vspace{0.6em}

    \begin{subfigure}[b]{0.48\textwidth}
        \centering
        \includegraphics[width=\linewidth]{fig3_training_loss.png}
        \caption{Training convergence.}
    \end{subfigure}
    \hfill
    \begin{subfigure}[b]{0.48\textwidth}
        \centering
        \includegraphics[width=\linewidth]{fig4_architecture_ablation.png}
        \caption{Architecture ablation (F1-score).}
    \end{subfigure}
    \caption{MLP-based terrain classifier evaluation, final configuration (150--80 hidden units). (a) Confusion matrix on the held-out test set ($n=120$). (b) ROC curve (AUC = 0.864). (c) Training-loss convergence across epochs. (d) F1-score comparison across the four tested hidden-layer architectures, motivating the 150--80 choice.}
    \label{fig:mlp_eval_grid}
\end{figure*}

The MLP was further benchmarked against alternative classifiers (Logistic Regression, SVM with RBF kernel, and Random Forest) under the same cross-validation protocol. Logistic Regression underperformed ($\text{F1} = 0.692$), reflecting the non-linear separability of the IMU features, while SVM-RBF ($\text{F1} = 0.857$) and Random Forest ($\text{F1} = 0.947$) achieved higher raw scores (Figure~\ref{fig:classifier_tradeoff}a).

\revblue{Per-sample inference latency (Figure~\ref{fig:classifier_tradeoff}b) explains why the MLP was nonetheless selected: it requires only} $1.64\,\mu\text{s}$ per sample, approximately $14\times$ faster than the SVM-RBF ($23.5\,\mu\text{s}$) and $25\times$ faster than the Random Forest ($41.0\,\mu\text{s}$)\revblue{, while Logistic Regression, at a comparable $0.54\,\mu\text{s}$, is markedly less accurate}. \revblue{Given the real-time constraints of the neural-arbitration pipeline and the prospect of deploying the classifier on a resource-constrained embedded microcontroller,} the MLP (150–80) was therefore selected as the configuration offering the best accuracy-to-latency trade-off for terrain classification, rather than the configuration with the highest raw classification score. Figure~\ref{fig:classifier_tradeoff}c additionally reports model size (parameters/tree nodes) for the same four classifiers\revblue{, relevant for embedded deployment}. The $1.64\,\mu\text{s}$ figure above was measured on a host PC (PyTorch)\revblue{, the external laptop described in Section~\ref{ssec:HardwareSpec}}; the same forward pass was additionally re-implemented in C and benchmarked on the target ESP32-S3 microcontroller ($N=1000$ repeated inferences), yielding a mean on-device latency of $8.81\,\text{ms}$ ($\text{SD}=2.67\,\mu\text{s}$) and a memory footprint of 24{,}140~B RAM (7.4\% of 320~KB) and 397{,}829~B flash (30.4\% of the board's 4~MB partition, almost entirely the $\sim$47k-parameter model itself). The on-device number is markedly higher than the host figure, consistent with a straightforward, non-vectorized scalar implementation reading weights from flash rather than from a host CPU's cache hierarchy.

\begin{figure}[htbp]
    \centering
    \includegraphics[width=\textwidth]{fig_classifier_tradeoff_fused.png}
    \caption{Classification/deployment trade-off across baseline classifiers (Logistic Regression, SVM-RBF, Random Forest, and MLP). (a) Accuracy, precision, recall, and F1-score\revblue{, mean of 5-fold stratified cross-validation on the IMU dataset ($n=600$ samples)}. (b) Per-sample inference latency. (c) Model size (parameters/tree nodes).}
    \label{fig:classifier_tradeoff}
\end{figure}


\subsubsection{Response to Variable Topographies}

In the real implementation (Figure~\ref{fig:RES_IMU_real}), images were captured every 21 seconds. Initially, the robot advances over a region with minimal irregularities ---less than 1.2cm--- which it overcomes during the first 13 seconds. At this point, a stagnation condition is detected that persists for the next 10 seconds. When this stagnation signal reaches its maximum sustained amplitude, the transition to the quadrupedal mode is triggered. In this state, the robot traverses the rough terrain and, around second 86, when the activity of unit $X_{0}$ ceases (Figure~\ref{fig:NEU_IMU_real}), the platform transitions back to the differential rolling mode.

\begin{figure}[t]
  \centering
  \begin{subfigure}[t]{0.48\textwidth}
    \centering
    \includegraphics[width=\linewidth]{E1.png}
    \caption{Entering the terrain, $\approx13$~s.}
    \label{fig:frames-a2}
  \end{subfigure}\hfill
  \begin{subfigure}[t]{0.48\textwidth}
    \centering
    \includegraphics[width=\linewidth]{E3.png}
    \caption{Quadrupedal transformation at stagnation peak.}
    \label{fig:frames-c2}
  \end{subfigure}

  \vspace{0.6em}

  \begin{subfigure}[t]{0.48\textwidth}
    \centering
    \includegraphics[width=\linewidth]{E5.png}
    \caption{Sustained quadrupedal crossing.}
    \label{fig:frames-e2}
  \end{subfigure}\hfill
  \begin{subfigure}[t]{0.48\textwidth}
    \centering
    \includegraphics[width=\linewidth]{E6.png}
    \caption{Recovery to rolling mode, $t\approx86$~s.}
    \label{fig:frames-f2}
  \end{subfigure}

  \caption{Four representative frames from the real-robot experiment facing an irregular terrain: (a)–(d) show the temporal evolution across the trial, detailed in the accompanying text.}
  \label{fig:RES_IMU_real}
\end{figure}

% NEU_IMU_real merged with NEU_IMU2_real below, after the second-scenario prose (see fig:neu_imu_real_pair).

% GRAPH_IMU_real merged with GRAPH_IMU2_real below, after the second-scenario prose (see fig:graph_imu_real_pair).

A similar behavior is observed in the experiment with the physical robot (Figure~\ref{fig:RES_IMU2_real}). However, in the real robot's gait-decision network, the inclined terrain is detected before second 46 (Figure~\ref{fig:NEU_IMU2_real}, Locomotion Module panel); consequently, a transient response emerges, increasing the activity of neuron \revblue{$X_{16}$}, which commands the robot to climb the ramp using the differential rolling mode; this mode remains active even after clearing the slope. \revblue{MLP outputs and stagnation state for both scenarios are shown in Figures~\ref{fig:GRAPH_IMU_real} and~\ref{fig:GRAPH_IMU2_real}.}

\begin{figure}[t]
    \centering
    \includegraphics[width=1.0\textwidth]{RES_IMU2.0_real.png}
    \caption{Trajectory in a terrain with elevation in real test. The poses were taken every 14 seconds. \revblue{Single trial ($N=1$).}}
    \label{fig:RES_IMU2_real}
\end{figure}

\begin{figure}[t]
  \centering
  \begin{subfigure}[t]{0.48\textwidth}
    \centering
    \includegraphics[width=\linewidth]{NEU_IMU_real.png}
    \caption{Irregular terrain. Locomotion decision module (top group); gait mode control module (bottom group).}
    \label{fig:NEU_IMU_real}
  \end{subfigure}\hfill
  \begin{subfigure}[t]{0.48\textwidth}
    \centering
    \includegraphics[width=\linewidth]{NEU_IMU2_real.png}
    \caption{Inclined terrain. Locomotion decision module (top group); gait mode control module (bottom group).}
    \label{fig:NEU_IMU2_real}
  \end{subfigure}
  \caption{Neural activations during locomotion over the two physical terrain scenarios: (a) irregular terrain; (b) inclined terrain. Within each panel, the Locomotion Module group appears above the Gait Decision Module group. \revblue{Each panel shows a single trial ($N=1$ per panel).}}
  \label{fig:neu_imu_real_pair}
\end{figure}

\begin{figure}[t]
  \centering
  \begin{subfigure}[t]{0.48\textwidth}
    \centering
    \includegraphics[width=\linewidth]{GRAPH_IMU_real.png}
    \caption{MLP outputs and IMU state timeline on uneven terrain.}
    \label{fig:GRAPH_IMU_real}
  \end{subfigure}\hfill
  \begin{subfigure}[t]{0.48\textwidth}
    \centering
    \includegraphics[width=\linewidth]{GRAPH_IMU2_real.png}
    \caption{MLP outputs and IMU state timeline on an inclined terrain.}
    \label{fig:GRAPH_IMU2_real}
  \end{subfigure}
  \caption{MLP-output and IMU-state timelines motivating the mode transition in each physical terrain scenario: (a) uneven terrain; (b) inclined terrain. \revblue{Each panel shows a single trial ($N=1$ per panel).}}
  \label{fig:graph_imu_real_pair}
\end{figure}

To further consolidate the physical validation under complex, multi-section topographies, a dedicated testing suite of 15 continuous repetitive trials was conducted over an experimental track whose structural layout systematically combines four topographic segments: a double-ascent ramp, a horizontal elevated plane, a controlled descent ramp, and a rough terrain section characterized by stochastic, non-uniform height variations generated via Perlin noise. \revblue{Figure~\ref{fig:real_variable_topography_sequence} and Table~\ref{tab:terrain_sequence} document the resulting gait-switching sequence.}

\begin{table}[htbp]
\centering
\caption{Gait-mode transitions across the six-frame compound-terrain sequence (Figure~\ref{fig:real_variable_topography_sequence}), replicated across 15 continuous trials.}
\label{tab:terrain_sequence}
\begin{tabularx}{\textwidth}{>{\centering\arraybackslash}p{0.08\textwidth} >{\raggedright\arraybackslash}p{0.19\textwidth} >{\raggedright\arraybackslash}p{0.19\textwidth} >{\raggedright\arraybackslash}X}
\toprule
\textbf{Frame} & \textbf{Terrain segment} & \textbf{Mode} & \textbf{Trigger / mechanism} \\
\midrule
1 & \revblue{Ramp entry (start)} & Quadrupedal Walking & \revblue{Default gait at the start; static stability and traction priority} \\
2 & \revblue{Double-ramp ascent} & \revblue{Quadrupedal $\rightarrow$ Differential Rolling} & \revblue{Incline too steep for the quadrupedal gait: pitch exceeds the threshold, suppressing quadrupedal salience and disinhibiting the rolling channel} \\
3 & \revblue{Horizontal plane (top of the ramp)} & Differential Rolling (sustained) & Energy-conservative wheeled traversal; rolling channel remains stable \\
4 & Controlled ramp descent & Differential Rolling (sustained) & Pitch disturbance stays below quadrupedal re-engagement threshold ($\approx 5^{\circ}$) \\
5 & Perlin-noise rough terrain & Rolling $\rightarrow$ Stagnation & Wheels lose traction on high-variance surface; MLP stagnation channel triggers sustained signal \\
6 & Rough terrain (recovery) & $\rightarrow$ Quadrupedal Walking & Stagnation triggers reactive disinhibition of quadrupedal channel; limbs redeploy, restoring support polygon \\
\bottomrule
\end{tabularx}
\end{table}

\revblue{The arbitration network runs through this sequence without operator input.} The critical event occurs in Frames 5--6: rolling-mode wheel slip on the stochastic rough segment triggers a sustained stagnation signal, resolved by reactively redeploying the quadrupedal gait, which restores the support polygon and completes traversal. This closed-loop reactive switching was observed across all 15 trials, demonstrating reliable, autonomous failure detection and gait recovery under stochastic, non-structured terrain conditions.

\begin{figure*}[t]
  \centering
  % ---- Row 1 ----
  \begin{subfigure}[t]{0.31\textwidth}
    \centering
    \includegraphics[width=\linewidth]{Prueba_Terreno_Parte_1.png}
    \caption{Frame 1 (see Table~\ref{tab:terrain_sequence}).}
    \label{fig:real_terrain_frame1}
  \end{subfigure}\hfill
  \begin{subfigure}[t]{0.31\textwidth}
    \centering
    \includegraphics[width=\linewidth]{Prueba_Terreno_Parte_2.png}
    \caption{Frame 2.}
    \label{fig:real_terrain_frame2}
  \end{subfigure}\hfill
  \begin{subfigure}[t]{0.31\textwidth}
    \centering
    \includegraphics[width=\linewidth]{Prueba_Terreno_Parte_3.png}
    \caption{Frame 3.}
    \label{fig:real_terrain_frame3}
  \end{subfigure}

  \vspace{0.6em}

  % ---- Row 2 ----
  \begin{subfigure}[t]{0.31\textwidth}
    \centering
    \includegraphics[width=\linewidth]{Prueba_Terreno_Parte_4.png}
    \caption{Frame 4.}
    \label{fig:real_terrain_frame4}
  \end{subfigure}\hfill
  \begin{subfigure}[t]{0.31\textwidth}
    \centering
    \includegraphics[width=\linewidth]{Prueba_Terreno_Parte_5.png}
    \caption{Frame 5.}
    \label{fig:real_terrain_frame5}
  \end{subfigure}\hfill
  \begin{subfigure}[t]{0.31\textwidth}
    \centering
    \includegraphics[width=\linewidth]{Prueba_Terreno_Parte_6.png}
    \caption{Frame 6.}
    \label{fig:real_terrain_frame6}
  \end{subfigure}

  \caption{Sequential 6-frame documentation of the physical platform navigating the compound variable-topography course, replicated across 15 continuous trials; per-frame terrain segment, mode, and trigger mechanism are detailed in Table~\ref{tab:terrain_sequence}.}
  \label{fig:real_variable_topography_sequence}
\end{figure*}

\subsection{Quantitative Performance Analysis in Physical Multi-Stimuli and Variable Terrain Environments}
\label{ssec:QuantReal}

\revblue{This subsection reports four stimulus scenarios on the physical platform (three hardware replicates each, consolidated in Table~\ref{tab:consolidated_physical_metrics}), followed by energy and latency profiling over variable terrain.}

%% ─── APPETITIVE SCENARIO ───────────────────────────────────────────────────

\begin{figure*}[t]
  \centering
  \begin{subfigure}[b]{0.48\textwidth}
    \centering
    \includegraphics[width=\linewidth]{Appetitive_Real_Plot_1_Switching_Delay_ECDF.png}
    \caption{Switching delay ECDF.}
    \label{fig:appetitive_ecdf_real}
  \end{subfigure}
  \hfill
  \begin{subfigure}[b]{0.48\textwidth}
    \centering
    \includegraphics[width=\linewidth]{Appetitive_Real_Plot_3_Temporal_Dynamics_PitchRollRMS.png}
    \caption{Pitch/Roll RMS attitude dynamics.}
    \label{fig:appetitive_rms_real}
  \end{subfigure}
  \caption{Real-world Appetitive Stimulus Scenario (\textit{Appetitive\_Real} family, three hardware replicates). (a) ECDF of locomotion mode-switching latency $T_{switch}$. (b) Coupled neural and kinematic response: (top) temporal evolution of the basal ganglia firing variance $Var(z)$ with red/blue stimulus-presence shading; (bottom) Roll and Pitch RMS attitude profile throughout the trial.}
  \label{fig:appetitive_real_compound}
\end{figure*}

\textbf{Appetitive (see Figure~\ref{fig:appetitive_real_compound}):} \revblue{The extended switching-delay window comes mainly from the electro-mechanical trajectory smoothing and landing interpolation enforced by the low-level controllers during target approach, which prevent peak-torque damage to the servo actuators.} The prolonged pre-lock exploration phase reflects real-world camera field-of-view and illumination constraints; the sustained, non-zero firing variance after mission completion reflects that the stopping neuron ($X_{17}$) freezes the physical stance while upstream sensory pathways continue registering the stimulus, rather than a full sensory shutdown. The cyclic Roll/Pitch oscillations stem from body-shaking during the rhythmic weight-shifting quadrupedal gait over uncalibrated floor textures, flattening once the platform converges to a stable equilibrium.

%% ─── AVERSIVE SCENARIO ────────────────────────────────────────────────────

\begin{figure*}[t]
  \centering
  \begin{subfigure}[b]{0.48\textwidth}
    \centering
    \includegraphics[width=\linewidth]{Aversive_Real_Plot_1_Switching_Delay_ECDF.png}
    \caption{Switching delay ECDF.}
    \label{fig:aversive_ecdf_real}
  \end{subfigure}
  \hfill
  \begin{subfigure}[b]{0.48\textwidth}
    \centering
    \includegraphics[width=\linewidth]{Aversive_Real_Plot_3_Temporal_Dynamics_PitchRollRMS.png}
    \caption{Pitch/Roll RMS attitude dynamics.}
    \label{fig:aversive_rms_real}
  \end{subfigure}
  \caption{Real-world Aversive Stimulus Scenario (\textit{Aversive\_Real} family, three hardware replicates). (a) ECDF of mode-switching latency $T_{switch}$ revealing the near-instantaneous emergency reflex. (b) Coupled neural and kinematic response: (top) neural firing variance $Var(z)$ displaying the dual-peak evasion signature, with red/blue stimulus-presence shading; (bottom) chassis attitude micro-dynamics during and after the evasive maneuver.}
  \label{fig:aversive_real_compound}
\end{figure*}

\textbf{Aversive (see Figure~\ref{fig:aversive_real_compound}):} Unlike the Appetitive scenario's extended settling window, the Aversive pipeline behaves as an immediate emergency reflex: upon threat detection, the ROS~2 Control architecture bypasses heavy trajectory filtering to prioritize high-speed backward flight, minimizing routing overhead to near-zero computational ticks under critical hazards. The dual-peak firing-variance pattern reflects, first, detection of the aversive stimulus and dispatch of the mode-switching command, then a second, higher-energy peak as hardware inertia keeps the threat briefly within the camera's field of view during the retreat, decaying back to its resting state once the chassis clears the threat from view. The concurrent Pitch step-increase and cyclic Roll oscillations reflect transient longitudinal tilt during the wheel-driven evasion, returning to a flat baseline once the threat leaves the visual field.

%% ─── COMPLEX MULTI-STIMULUS SCENARIO ─────────────────────────────────────

\begin{figure*}[t]
  \centering
  \begin{subfigure}[b]{0.48\textwidth}
    \centering
    \includegraphics[width=\linewidth]{Complex_Real_Plot_1_Switching_Delay_ECDF.png}
    \caption{Switching delay ECDF.}
    \label{fig:complex_ecdf_real}
  \end{subfigure}
  \hfill
  \begin{subfigure}[b]{0.48\textwidth}
    \centering
    \includegraphics[width=\linewidth]{Complex_Real_Plot_3_Temporal_Dynamics_PitchRollRMS.png}
    \caption{Pitch/Roll RMS attitude dynamics.}
    \label{fig:complex_rms_real}
  \end{subfigure}
  \caption{Real-world Complex Multi-Stimulus Scenario (\textit{Complex\_Real} family, three hardware replicates). (a) ECDF of mode-switching latency under chained sequential transformations. (b) Coupled neural and kinematic response: (top) three-phase neural variance profile mapping obstacle avoidance, aversive escape, and appetitive convergence phases, with red/blue stimulus-presence shading; (bottom) step-like attitude excursions during transitions and final stabilization at ${\approx}\,4^{\circ}$.}
  \label{fig:complex_real_compound}
\end{figure*}

\textbf{Complex (see Figure~\ref{fig:complex_real_compound}):} \revblue{The step response stays deterministic across back-to-back transformations: the ROS~2 Control layer absorbs the electromechanical inertia and blends the trajectories, so torque stress does not accumulate in the joints.} The three-phase firing-variance profile maps the multi-stage decision process: a transient hyper-excitation peak during static obstacle avoidance, \revblue{driven by the dense LiDAR returns from the physical barrier}; a secondary activation wave capturing the emergency flight reflex to the aversive cue; and a low-energy steady-state as the robot tracks the appetitive target to mission completion. The chassis attitude telemetry mirrors this sequence, with asymmetric flanking maneuvers exciting the transversal axis and rapid backward accelerations peaking the longitudinal tilt, before both metrics drop and stabilize during the final convergence phase, indicating the wheel-driven configuration smooths parasitic oscillations during target approach.

%% ─── STATIC OBSTACLE AVOIDANCE SCENARIO ──────────────────────────────────

\begin{figure*}[t]
  \centering
  \begin{subfigure}[b]{0.48\textwidth}
    \centering
    \includegraphics[width=\linewidth]{Obstacle_Real_Plot_1_Switching_Delay_ECDF.png}
    \caption{Switching delay ECDF.}
    \label{fig:obstacle_ecdf_real}
  \end{subfigure}
  \hfill
  \begin{subfigure}[b]{0.48\textwidth}
    \centering
    \includegraphics[width=\linewidth]{Obstacle_Real_Plot_3_Temporal_Dynamics_PitchRollRMS.png}
    \caption{Pitch/Roll RMS attitude dynamics.}
    \label{fig:obstacle_rms_real}
  \end{subfigure}
  \caption{Real-world Static Obstacle Avoidance Scenario (\textit{Obstacle\_Real} family, three hardware replicates). (a) Instantaneous ECDF of mode-switching latency. (b) Coupled neural and kinematic response: (top) massive transient hyper-excitation peak in neural firing variance $Var(z)$ induced by volumetric geometric input (no stimulus-presence shading---this scenario uses a non-chromatic obstacle cue); (bottom) rhythmic Roll RMS expansion during active lateral flanking and post-avoidance stabilization.}
  \label{fig:obstacle_real_compound}
\end{figure*}

\textbf{Obstacle (see Figure~\ref{fig:obstacle_real_compound}):} \revblue{As in the Aversive scenario, the step response is immediate and consistent across all three replicates: geometric blockages are handled as emergency reflexes without trajectory ramp filtering.} The massive transient spike in firing variance---an order of magnitude above the color-cue scenarios---\revblue{reflects the obstacle entering the LiDAR detection range and activating its obstacle-encoding ring}, triggering a hyper-excitation wave that drives aggressive lateral flanking, collapsing back to baseline once the chassis clears the obstacle's spatial boundary. The rhythmic Roll expansion with coupled Pitch variation reflects the weight-shifting needed for a tight lateral bypass, flattening once safety is secured.

\begin{table}[htbp]
\centering
\caption{Consolidated Quantitative Metrics for the Physical Multi-Stimuli Scenarios (three hardware replicates each). $T_{switch}$: mode-switching-delay ECDF window. $Var(z)$: basal-ganglia firing-variance peak/pattern. Roll/Pitch RMS: chassis attitude range during the active response window.}
\label{tab:consolidated_physical_metrics}
\begin{tabularx}{\textwidth}{l >{\centering\arraybackslash}X >{\centering\arraybackslash}X >{\centering\arraybackslash}X >{\centering\arraybackslash}X}
\toprule
\textbf{Scenario} & \textbf{$T_{switch}$ ECDF} & \textbf{$Var(z)$ peak} & \textbf{Roll / Pitch RMS} & \textbf{Settles by} \\
\midrule
Appetitive & $3.0$--$9.0\,\text{s}$ (mechanically limited) & Sustained, non-zero post-completion & $3.0^{\circ}$--$3.5^{\circ}$ (cyclic) & $t > 25\,\text{s}$ \\
Aversive   & ${\approx}\,0.0\,\text{s}$ (instantaneous) & Dual-peak (detection + retreat-FOV) & Pitch step $\uparrow$, Roll cyclic ($t{=}4$--$9\,\text{s}$) & $t > 10\,\text{s}$ \\
Complex    & ${\approx}\,1.0$--$1.2\,\text{s}$ & Three-phase, peak ${\approx}\,27\times$ Obstacle baseline & Roll ${\approx}\,7^{\circ}$, Pitch ${\approx}\,5^{\circ} \to 4^{\circ}$ & $t > 25\,\text{s}$ \\
Obstacle   & ${\approx}\,0.0\,\text{s}$ (instantaneous) & ${\approx}\,1100$, transient ($t < 7\,\text{s}$) & Roll $3.5^{\circ}$ (rhythmic) & $t = 8\,\text{s}$ \\
\bottomrule
\end{tabularx}
\end{table}

%% ─── ENERGY CONSUMPTION ───────────────────────────────────────────────────

\begin{figure}[t]
  \centering
  \includegraphics[width=0.5\textwidth]{FigReal_Morphological_Transition_Energy.png}
  \caption{Morphological transition energy per direction, derived from $|I_{\text{current}}| \times 7.2\,\text{V}$ \revblue{(current sampled at $\approx 5$\,Hz by the onboard INA sensor; fixed nominal voltage of the NiMH battery, Table~\ref{table:HardwareParams})} integrated over the pre-calibrated electromechanical coupling window (2.92\,s / 1.56\,s, matching Figure~\ref{fig:FourPanelLatencies}), across $N=16$ (Quadrupedal $\rightarrow$ Differential Rolling) and $N=19$ (Differential Rolling $\rightarrow$ Quadrupedal) individually detected structural-command transitions. Error bars: $\pm 1$ SD; the between-direction difference is not statistically significant (Welch's $t \approx 1.18$).}
  \label{fig:MorphologicalEnergy}
\end{figure}

The energetic cost of physical adaptation was quantified from individually detected structural-command transitions \revblue{as described in Figure~\ref{fig:MorphologicalEnergy}'s caption}. The Quadrupedal $\rightarrow$ Differential Rolling transition required $5.06 \pm 2.93$~J ($N=16$), and the reverse transition (Differential Rolling $\rightarrow$ Quadrupedal) required $4.13 \pm 1.19$~J ($N=19$); the between-direction difference is not statistically significant \revblue{(see caption)}. Greater trial-to-trial variability was observed in the Quadrupedal $\rightarrow$ Differential Rolling direction, though the present dataset does not isolate a specific mechanical cause for this dispersion. As the Quadrupedal$\leftrightarrow$Differential Rolling pair involves the largest morphological reconfiguration among the three modes (full limb deployment/retraction together with the leg-to-wheel locomotion change), its transition energy serves as an upper bound on the cost of any mode transition: in the worst case a transition consumes $\leq 5.06$~J, a small cost relative to continuous operation (e.g.\ a typical 38\,s Omnidirectional segment integrates to $\approx 216$~J). \revblue{Table~\ref{tab:mode_steady_state_energy} reports the long-term average steady-state power draw of each mode.}

\begin{table}[htbp]
\centering
\caption{Long-term average steady-state power consumption per locomotion mode, computed from
pooled intra-segment power samples (individual samples within each structural-command-detected mode
segment, electromechanical coupling window excluded; see Figure~\ref{fig:MorphologicalEnergy}). Power
derived as $|I|\times7.2\,\text{V}$\revblue{, with current sampled at $\approx 5$\,Hz by the onboard INA sensor and 7.2\,V the nominal voltage of the 3300\,mAh NiMH battery (Table~\ref{table:HardwareParams}). Mean and SD are taken over all individual samples pooled across the segments of each mode.}}
\label{tab:mode_steady_state_energy}
\begin{tabularx}{\textwidth}{l >{\centering\arraybackslash}X >{\centering\arraybackslash}X}
\toprule
\textbf{Locomotion Mode} & \textbf{N (segments)} & \textbf{Mean Power $\pm$ SD (W)} \\
\midrule
Quadrupedal ($C$) & 20 & $2.73 \pm 2.29$ \\
Differential Rolling ($H$) & 20 & $1.13 \pm 1.08$ \\
Omnidirectional ($X$) & 1 & $5.71 \pm 0.73$ \\
\bottomrule
\end{tabularx}
\\[2pt]
\footnotesize{Intra-segment sample counts: $C$ $N=1{,}873$; $H$ $N=3{,}297$; $X$ $N=194$ (single
38.6\,s segment; hardware subsequently decommissioned, no further Omnidirectional recordings
possible).}
\end{table}

This roughly five-fold gap between Omnidirectional and Differential Rolling has a direct mechanical explanation. In Differential Rolling, wheel-floor friction is aligned with the direction of travel and acts as useful propulsive force; in Omnidirectional mode, the diagonal wheel configuration must instead continuously overcome lateral friction that contributes nothing to displacement and is dissipated as heat, requiring sustained extra motor power. Quadrupedal walking falls in between, its cost dominated by ground contact/lift-off impact losses rather than continuous lateral friction. A mission controller should accordingly minimize time spent in Omnidirectional mode, reserving it for maneuvers that strictly require omnidirectional mobility.

%% ─── TERRAIN LATENCY (4 PANELS) ───────────────────────────────────────────

\begin{figure}[t]
  \centering
  \includegraphics[width=0.90\textwidth]{FigReal_Terrain_Latencies_4Panels.png}
  \caption{Four-panel latency partitioning for real-world autonomous multi-terrain navigation (15 hardware replicates). Panel~A: neural processing latency, Quadrupedal $\rightarrow$ Differential Rolling ($\mu = 0.184\,\text{s}$). Panel~B: neural processing latency, Differential Rolling $\rightarrow$ Quadrupedal ($\mu = 0.534\,\text{s}$). Panel~C: electromechanical coupling latency, Quadrupedal $\rightarrow$ Differential Rolling ($\mu = 2.92\,\text{s}$). Panel~D: emergency retraction latency, Differential Rolling $\rightarrow$ Quadrupedal ($\mu = 1.56\,\text{s}$).}
  \label{fig:FourPanelLatencies}
\end{figure}

To evaluate repeatability and determinism under uncalibrated physical environments, a continuous multi-terrain navigation mission was executed across 15 hardware replicates. The track forced the robot from Quadrupedal to Differential Rolling Mode to climb a steep incline, then back to Quadrupedal Mode on unstructured terrain inducing wheel slippage; the resulting temporal dynamics are partitioned into a four-panel layout to preserve resolution against low statistical variance (Figure~\ref{fig:FourPanelLatencies}). The basal ganglia loop's decision latency shows a clear task-dependent split: the Quadrupedal $\rightarrow$ Differential Rolling transition (Panel~A) achieves an instantaneous step response, reflecting immediate pitch-gradient adaptation, while the reverse transition (Panel~B), triggered by wheel slippage on rugged terrain, settles at a markedly longer latency --- a deliberate stability feature that filters transient high-frequency impact noise to suppress control chattering and false-positive switches on noisy surfaces. The electromechanical coupling timelines reflect these commands mechanically: achieving full wheel-ground alignment (Panel~C) requires an extended settling window, capturing the kinematic cost of leg extension under slope conditions, whereas the emergency retraction to the high-traction articulated stance (Panel~D) is a rapid recovery reflex, minimizing the wheel-spin window and securing foothold anchors quickly on unstable terrain.

\subsection{Limitations of the Present Work}
\label{ssec:Limitations}

As an open, low-cost exploratory research baseline, the physical robot is constructed using low-torque servomotors and a lightweight PLA structure, which limits its operational envelope compared to industrial-grade platforms. Locomotion gaits are currently executed via pre-configured open-loop control patterns; while this reduces control complexity, it weakens active stabilization capabilities on high-slip surfaces or excessively steep inclinations. \revblue{This limit lies in the hardware and is independent of the neural arbitration architecture.} As a representative case, the gait selected during steep-slope negotiation (Figure~\ref{fig:RES_IMU2_real}) reflects the boundaries of the physical actuator envelope: the neural decision module consistently selects the locomotion mode associated with improved stability under the detected terrain conditions, yet the open-loop execution of that gait cannot compensate for the torque deficits inherent to the hardware. Experimental validation was constrained to structured laboratory environments with static stimuli and did not include unconstrained outdoor topographies, dynamic obstacles, or severe external disturbances. Therefore, the empirical results presented herein do not constitute validation of a production-ready system for rugged industrial deployment. Instead, they substantiate the mathematical, biological, and computational feasibility of a highly efficient basal-ganglia inspired arbitration network capable of managing multimodal mobility \revblue{on a low-cost robotic platform}.


\revblue{Regarding the computing platform: the neural circuits, the MLP terrain classifier, and the image processing ran on the external laptop of Section~\ref{ssec:HardwareSpec}, not on the robot's ESP32-S3. Onboard feasibility was assessed only for the MLP classifier (Section~\ref{sssec:mlp_eval}); the basal-ganglia, locomotion, and gait-decision circuits and the vision pipeline were not benchmarked on the microcontroller, so a fully autonomous onboard implementation remains future work.}

%
Regarding the virtual test environment: the Gazebo world used for
obstacle-evasion and stimulus-response experiments
(Figure~\ref{fig:depot_simulation}) is a symbolic representation
of an industrial storage facility.
It was not constructed from engineering measurements of any
real facility, and its spatial statistics (corridor width,
obstacle density, scene dimensions) have not been benchmarked
against actual infrastructure-inspection deployments.
Consequently, the simulation results should be interpreted as
demonstrations of algorithmic behavior in a \emph{representative}
constrained environment, rather than as performance guarantees
for deployment in specific industrial buildings.%


Beyond the hardware and simulation-environment constraints above, the neural and locomotion-control parameters were manually tuned and kept constant throughout all experiments, without applying adaptive optimization processes that could enhance performance in more diverse scenarios; system resilience was also not assessed under partial hardware failures or temporary sensory dropout conditions likely to occur in real industrial environments \revblue{(beyond the simulated camera-detection dropout of Section~\ref{sssec:noise_robustness_extra})}. The multimodal sensory integration, moreover, does not follow a formal fusion framework in the classical signal-processing sense---such as Kalman filtering, Bayesian estimation, or deep sensor fusion. Instead, cross-modal integration emerges implicitly within the bio-inspired neural network: each sensory channel feeds differentiated neural pathways that converge through winner-take-all competition dynamics in the basal ganglia module, effectively distributing the fusion function across the network topology. While this approach is biologically grounded, it has not been formally compared against standard fusion methods. Similarly, visual stimulus detection relies on color-based segmentation calibrated under controlled illumination, without evaluation under varying lighting conditions or occlusion. Notwithstanding these design simplifications, the perception-decision noise-robustness evaluation (Section~\ref{ssec:NoiseRobustnessEval}) demonstrated empirical robustness under combined IMU-drift, LiDAR-noise, and camera-illumination-distortion perturbation, evaluated jointly rather than isolated per sensor --- no configuration in this evaluation exercises a single sensor's degradation in isolation from the other two. Within the same evaluation, additive IMU white noise (distinct from the random-walk drift that was exercised) was held at zero throughout and so was not tested here. Robustness was also not uniform across scenarios: the Rugged Terrain suite showed a markedly elevated failure rate under perturbation relative to the other five suites, with 8 tipovers and 4 crashes recorded out of 27 attempts (Section~\ref{sssec:noise_robustness_extra}), indicating that combined sensory perturbation degrades structural stability specifically on this terrain even though the trials that do complete show no corresponding trend toward instability. These constraints do not invalidate the results but indicate that the next development stage should focus on validating the proposed architecture in a prototype that includes robustness and adaptability tests under real, non-controlled conditions.

Regarding the basal-ganglia ablation study (Section~\ref{sec:results}): the conflict-efficiency metric $\lambda$ was excluded from the reported results because its current implementation only registers active conflict when two specific channels are simultaneously present within the sampling window, a condition rarely met in practice across these trials --- this is a limitation of the metric's present instrumentation, not evidence that conflict efficiency is unaffected by the ablations. Additionally, per-cycle firing-variance and temporal-consistency diagnostics, which could further characterize the \emph{threshold-only} variant's failure mode (e.g.\ oscillation near the decision boundary, as opposed to outright deadlock), were not captured for this dataset; the threshold-versus-WTA comparison therefore relies on decision latency and attitude stability rather than on these more targeted diagnostics.

The Reactive Inspection Task's scope constraints ---absence of deliberate path planning or mapping, the resulting inapplicability of a formal coverage-rate metric, and the first- rather than re-acquisition semantics of $T_{acquisition}$--- are established alongside that experiment's own results in Section~\ref{sssec:reactive_inspection} above and are not restated here.

